\documentclass[10pt,a4paper]{article}

\usepackage[utf8]{inputenc}
\usepackage[T1]{fontenc}
\usepackage[english]{babel}
\usepackage[margin=2.3cm,headheight=14pt]{geometry}
\usepackage{microtype}
\usepackage{amsmath,amssymb,mathtools}
\usepackage{booktabs,tabularx}
\usepackage{multirow}
\usepackage{longtable}
\usepackage{graphicx}
\usepackage{caption}
\usepackage{float}
\usepackage{placeins}
\usepackage{xcolor}
\usepackage{fancyhdr}
\usepackage[hidelinks]{hyperref}
\usepackage[numbers,sort&compress]{natbib}
\usepackage{algorithm}
\usepackage{algpseudocode}

\pagestyle{fancy}
\fancyhf{}
\fancyhead[L]{\footnotesize\itshape Trajectory-aware allocation of a
spherical-harmonic degree budget}
\fancyhead[R]{\footnotesize\thepage}
\renewcommand{\headrulewidth}{0.3pt}
\setlength{\parskip}{0pt}
\setlength{\parindent}{1.2em}
\captionsetup{font=small,labelfont=bf,labelsep=period,skip=4pt}

\renewcommand{\arraystretch}{1.12}
\setlength{\tabcolsep}{5pt}
\setlength{\heavyrulewidth}{0.9pt}
\setlength{\lightrulewidth}{0.5pt}
\setlength{\aboverulesep}{2.2pt}
\setlength{\belowrulesep}{2.2pt}

\newcommand{\papertitle}{Trajectory-aware allocation of a spherical-harmonic
degree budget: a non-separable formulation, an attainability benchmark, and a
deployable controller}

\hypersetup{
  pdftitle={\papertitle},
  pdfauthor={Ayberk Demirkanat},
  pdfsubject={Trajectory-aware spherical-harmonic degree allocation under a
    declared compute budget},
  pdfkeywords={lunar gravity field; spherical harmonics; truncation degree;
    resource allocation; state-transition matrix; orbit propagation; GRAIL}
}

\renewcommand{\topfraction}{0.92}
\renewcommand{\bottomfraction}{0.80}
\renewcommand{\textfraction}{0.06}
\renewcommand{\floatpagefraction}{0.80}
\setcounter{topnumber}{3}
\setcounter{bottomnumber}{2}
\setcounter{totalnumber}{4}

\newcommand{\code}[1]{\texttt{#1}}

% ===========================================================================
% DRAFT MACHINERY -- remove before submission
% ===========================================================================
%
% \ph{...}   numeric placeholder. Every one of these must eventually be
%            replaced by a macro generated into metrics/ by the campaign
%            scripts, exactly as the previous paper did, so that no number is
%            typed by hand into the prose. Find them all with:
%                grep -rn '\\ph{' chapters/
%
% \dnote{..} a note about which way the text will have to go depending on the
%            campaign outcome. These mark the paragraphs whose DIRECTION is
%            not yet known -- they are not editorial to-dos.
%
% \wpref{..} the work package (see ../WP.md) that produces the numbers for the
%            surrounding passage.
%
% Set \draftfalse to typeset a clean version for reading; placeholders then
% render as a plain marker rather than in colour, and \dnote disappears.

\newif\ifdraft
\drafttrue

% \ph is \ensuremath-wrapped so that it is legal inside and outside math mode.
\ifdraft
  \newcommand{\ph}[1]{\ensuremath{\langle\,\text{\color{red}\sffamily #1}\,\rangle}}
  \newcommand{\dnote}[1]{\par\vspace{2pt}\noindent
    {\small\color{blue!60!black}\textsf{[direction] #1}}\par\vspace{2pt}}
  \newcommand{\wpref}[1]{\textcolor{green!35!black}{\textsf{\footnotesize[#1]}}}
\else
  \newcommand{\ph}[1]{\textbf{??}}
  \newcommand{\dnote}[1]{}
  \newcommand{\wpref}[1]{}
\fi

% Shorthand for the policy codes of NOTATION.md, so a rename is one edit.
\newcommand{\Fop}{\textsc{f-op}}
\newcommand{\Fenv}{\textsc{f-env}}
\newcommand{\Rrad}{\textsc{r-rad}}
\newcommand{\Rint}{\textsc{r-int}}
\newcommand{\Aforce}{\textsc{a-force}}
\newcommand{\Asens}{\textsc{a-sens}}
\newcommand{\Asign}{\textsc{a-sign}}
\newcommand{\Cplan}{\textsc{c-plan}}
\newcommand{\Cfb}{\textsc{c-fb}}
\newcommand{\Ctgo}{\textsc{c-tgo}}
\newcommand{\Clite}{\textsc{c-lite}}

\DeclareMathOperator*{\argmin}{arg\,min}
\newcommand{\dd}{\mathrm{d}}


\begin{document}

\thispagestyle{plain}

\begin{center}
  {\LARGE\bfseries \papertitle\par}
  \vspace{14pt}
  {\large Ayberk Demirkanat\par}
  \vspace{4pt}
  {\small\itshape Department of Astronautical Engineering, Istanbul
  Technical University, Istanbul, T\"urkiye\par}
  \vspace{2pt}
  {\small \code{demirkanat21@itu.edu.tr}\par}
\end{center}

\vspace{10pt}
\begin{center}
\begin{minipage}{0.92\textwidth}
  \small
  \textbf{Abstract.}
  A spherical-harmonic evaluation costs $\mathcal{O}(N^2)$, so choosing the
  truncation degree along an orbit is a resource-allocation problem. Allocating
  from the instantaneous force defect is known to order the arcs wrongly: a
  trajectory error is a signed integral of the defect against the
  state-transition matrix, and a defect norm is a sum of magnitudes. Writing
  the objective in the form the dynamics impose makes the allocation problem
  \emph{non-separable} in the per-epoch degrees, which is why the standard
  Lagrangian treatment of a budgeted separable allocation does not apply. We
  show that the discretized objective is a quadratic form whose coupling matrix
  depends on the epoch pair only through $\max(i,k)$, so a full sweep over the
  $K_{\mathrm{dec}}$ decision intervals, evaluated from $M$ accumulation
  samples over $\lvert\mathcal{N}\rvert$ candidate degrees, costs
  $\mathcal{O}(M\lvert\mathcal{N}\rvert)$, and a Frank--Wolfe bound on the
  convex-hull relaxation certifies the gap; the resulting allocation is
  obtained without propagating anything, and every policy compared switches on
  one common decision grid so that the benchmark's margin is information rather
  than resolution. On \ph{n} orbits of \ph{n} independent coverage designs
  the certified benchmark is a median factor of \ph{value} better than the
  equal-budget constant degree. The three criteria the paper separates are
  three readings of one trajectory kernel --- ignore it, retain its local
  diagonal, retain the full cross-epoch coupling --- so the question of what
  sensitivity weighting misses becomes the question of how much lives off the
  diagonal, and the answer is \ph{pct}.
  Deploying this requires the \emph{direction} of the omitted acceleration, not
  only its magnitude, and that direction varies on the field's own resolution
  scale $\pi r/N$, so a surrogate coarser than that scale cannot represent it
  and one finer carries as many degrees of freedom as the omitted coefficients.
  We estimate it instead from the leading omitted degree bands, probing forward
  along the vehicle's own short-horizon prediction, and complete its amplitude
  from the model's measured degree-variance spectrum. That split is the paper's
  information statement: the spectrum is one-dimensional and costs nothing to
  upload, while the texture is a field at resolution $\pi r/N$ and has to be
  measured in flight. Covering a decision interval costs a measured \ph{pct} of
  the gravity budget --- the controller's dominant overhead, charged to it in
  full. The
  resulting controller plans a cancellation target offline on a cheap pilot arc
  through a forward--backward iteration and probes forward, along its own
  short-horizon state prediction, before committing each decision interval. Held to equal \emph{realized}
  quadratic work including its own overhead, it \ph{outcome} the constant
  degree on \ph{n} of \ph{n} resolved comparisons across \ph{n} populations.
  All conclusions are model-relative and bounded by the tested designs.

  \vspace{6pt}
  \textbf{Keywords:} lunar gravity field; spherical harmonics; truncation
  degree; degree allocation; state-transition matrix; orbit propagation; GRAIL
\end{minipage}
\end{center}
\vspace{6pt}

\dnote{The abstract is written to the \textcolor{blue!60!black}{green} outcome
band of \code{../OUTCOMES.md}. Three of its clauses invert if the campaign
lands in the amber or orange band, and the alternative wordings are drafted at
the head of \code{chapters/11\_conclusion.tex} so that the paper is not
rewritten under time pressure.}

\section*{Nomenclature}
\noindent
\begin{minipage}[t]{0.48\textwidth}
\emph{Symbols}\\[2pt]
\begin{tabular}{@{}l@{\ \ }l@{}}
$N$ & spherical-harmonic truncation degree \\
$N_i$ & degree in force at accumulation epoch $i$ \\
$\mathcal{N}$ & candidate degree set \\
$N_{\mathrm{crit}}$ & critical-altitude fixed degree \\
$R$ & reference radius of the field model \\
$M$ & number of accumulation epochs \\
$\Delta t_{\mathrm{acc}}$ & accumulation step \\
$\Delta t_{\mathrm{dec}}$ & decision step (degree held constant) \\
$\Delta t_{\mathrm{probe}}$ & tail-probe refresh interval \\
$\Delta\mathbf a$ & truncation acceleration defect \\
$\hat{\mathbf v}$ & estimated omitted-acceleration vector \\
$\Phi$ & reference state-transition matrix \\
$\mathbf B$ & acceleration-to-velocity input map \\
$\mathbf G$ & gravity gradient $\partial\mathbf a/\partial\mathbf r$; generates $\Phi$ \\
$\mathbf u_i$ & epoch contribution $\Phi(t_0,t_i)\mathbf B\Delta\mathbf a\,\Delta t$ \\
$\mathbf S_j$ & partial sum $\sum_{i\le j}\mathbf u_i$ \\
$\mathbf Q$ & objective coupling matrix \\
$\mathbf c_i$ & cross term $\sum_k \mathbf Q_{ik}\mathbf u_k$ \\
$\lambda$ & budget multiplier \\
$\mathbf H_r$ & position selector $[\,\mathbf I_3\;\;\mathbf 0\,]$ \\
$\Delta t_i$ & accumulation step at epoch $i$ (non-uniform) \\
$\omega_j$ & outer-quadrature weight, $\omega_j\ge0$, $\sum_j\omega_j=T$ \\
$\mathbf K(\tau,\sigma)$, $\mathbf K_i$ & objective kernel and its diagonal $\mathbf K(t_i,t_i)$; the \Asens\ weight \\
$g(i)$, $I_q$ & epoch-to-interval map and its fibre \\
$W_q$ & interval time weight $\sum_{i\in I_q}\Delta t_i$ \\
$K_{\mathrm{dec}}$ & number of decision intervals \\
$T$ & arc length \\
$k$ & probe depth (number of leading omitted bands) \\
$\delta$ & candidate-window half-width, in grid indices \\
$\Delta_{\mathrm{span}}$ & degree span the window induces \\
$n_{\mathrm{probe}}$ & probe points per decision interval \\
$r_{\mathrm{RHS}}$, $N_{\mathrm{RHS}}$ & right-hand-side call rate and arc total \\
$\gamma$ & amplitude completion factor, from $P_n$ \\
$P_n$ & coefficient degree variance \\
$J$ & forward--backward iteration count \\
$L_{\mathrm{FW}}$ & Frank--Wolfe lower bound \\
$g_J$, $g_E$ & certified gap in objective and error \\
\end{tabular}
\end{minipage}
\hfill
\begin{minipage}[t]{0.48\textwidth}
\emph{Budgets and metrics}\\[2pt]
\begin{tabular}{@{}l@{\ \ }l@{}}
$B_1$ & nominal work rate, time average $\langle N^2\rangle_t$ \\
$B_2$ & realized total quadratic work $\sum_{\mathrm{RHS}}N^2$ \\
$B_3$ & measured serial gravity-kernel time \\
$B_+$ & controller overhead \\
$B$ & total nominal work, $B=B_1T=\sum_q W_qN_q^2$ \\
$\beta$ & allocation target $B_1/N_{\mathrm{crit}}^2$ \\
$\beta_2$ & realized $B_2/(K_{\mathrm{ref}}N_{\mathrm{crit}}^2)$ \\
$E$ & propagated seven-day position RMS error \\
$\hat E$ & first-order predicted error, $\hat E=\sqrt{J}$ \\
$\rho,\ \hat\rho$ & error ratio, comparator over candidate \\
$f$ & capture fraction of the benchmark gap \\
$M_{\mathrm{res}}$ & resolution margin; resolved if $>1$ \\
$\kappa$ & probe direction accuracy $\cos\angle(\hat{\mathbf v},\Delta\mathbf a)$ \\
$\kappa_{\mathrm{eff}}$ & the same, averaged over a decision interval \\
$\tau_{\mathrm{corr}}$ & defect-direction correlation time $\pi r/(Nv)$ \\
$T_{\mathrm{coh}}$ & pilot-arc coherence horizon \\
\end{tabular}
\\[8pt]
\emph{Policies}\\[2pt]
\begin{tabular}{@{}l@{\ \ }l@{}}
\Fop & operational fixed comparator \\
\Fenv & fixed-family lower envelope (post hoc) \\
\Rrad & budget-calibrated radial endpoint \\
\Rint & interior member \\
\Aforce & force-level allocation benchmark \\
\Asens & sensitivity-weighted benchmark \\
\Asign & signed trajectory-aware benchmark \\
\Cplan & planned trajectory-aware controller \\
\Cfb & budget-feedback variant \\
\Ctgo & time-to-go-only ablation \\
\end{tabular}
\end{minipage}

\vspace{6pt}

\section{Introduction}
\label{sec:intro}

Evaluating a spherical-harmonic gravity field to degree $N$ costs
$\mathcal{O}(N^2)$, and an adaptive-step integrator makes tens of thousands of
force calls per propagated day. On an eccentric orbit the required degree is
not constant along the arc, so the choice of $N$ is not one number but a
schedule, and choosing that schedule under a declared compute budget is a
resource-allocation problem rather than an accuracy question.

The natural way to allocate is from the local truncation defect: spend degrees
where the omitted acceleration is largest. That is what an attenuation
threshold, a Kaula-weighted tail criterion, or a radius-driven rule all do in
different guises, and it is what a Lagrangian treatment of a budgeted
separable allocation formalizes \citep{everett1963generalized,
shoham1988efficient}. It also produces the wrong ranking. Held to one gravity
budget on two independent coverage designs, a schedule concentrated near
perilune lowers the trajectory-averaged truncation force defect by a median
factor of about five and still gives the larger seven-day position error
\citep[\ph{self-citation of the previous paper}]{atallah2022radial}. The
mechanism is not a sampling artefact: a first-order displacement is a
\emph{signed} integral of the defect against the state-transition matrix, and a
defect norm is a sum of magnitudes, so the two agree only when the defect is
incoherent, which on these arcs it is not.

\dnote{The self-citation above is the previous paper. Once it has a DOI,
replace the placeholder and re-check that every claim attributed to it here is
one that paper actually makes at the confidence level it makes it.}

That leaves an obvious question unanswered, and this paper answers it. If the
force-level criterion is the wrong objective, what does the right one give, how
much is it worth, and can any of it be obtained without the reference-field
access the benchmark is granted? Three questions organize the paper.

\begin{description}
  \item[RQ1] Under trajectory-aware information, how much accuracy is
    \emph{attainable} at a fixed budget, relative to the constant degree and to
    the best radius-driven schedules?
  \item[RQ2] How much of that comes from weighting each epoch by its
    sensitivity, and how much from the signed cancellation that a magnitude
    criterion discards?
  \item[RQ3] How much survives without the reference trajectory and without
    pointwise reference-field evaluations along it --- the controller keeping
    only global, offline model metadata such as the degree variances?
\end{description}

Three things make the answers non-obvious, and each is a contribution.

\paragraph{The objective is non-separable.} Writing the trajectory error in the
form the dynamics impose puts the norm \emph{after} the signed integral, which
couples every pair of epochs through cancellation. The standard multiplier
treatment therefore does not apply, and neither does the obvious repair of
reweighting the per-epoch defect by its sensitivity: that corrects the weight
while still summing magnitudes. Section~\ref{sec:problem} states the problem
and Section~\ref{sec:algorithm} shows that it is nonetheless cheap to solve
well, because the discretized coupling matrix depends on the epoch pair only
through $\max(i,k)$. A full coordinate sweep costs
$\mathcal{O}(M\lvert\mathcal{N}\rvert)$, and a Frank--Wolfe bound on the
convex-hull relaxation \citep{frank1956algorithm,jaggi2013revisiting}
certifies how far the result sits from the relaxed optimum. Nothing is
propagated to obtain it. The degree is a decision variable on a coarse decision
grid rather than on the fine grid the objective is accumulated over, which
costs nothing in the solve and matters for the comparison: a benchmark allowed
to switch faster than any deployable controller would earn part of its margin
from resolution rather than from information.

\paragraph{Deployment needs a direction, not a magnitude.} A compact tail rule
predicts how much acceleration a truncation omits; it says nothing about which
way it points. Cancellation is exactly the part that needs the direction, so a
controller built on a tail magnitude alone can reach the
sensitivity-weighted level and no further. Section~\ref{sec:controller} shows
that tabulating the direction is not a cheap way out either: the omitted field
varies on the model's own resolution scale $\pi r/N$, which at degree 300 is
under twenty kilometres of ground track, so a surrogate coarser than that scale
cannot represent it and one fine enough to do so carries a comparable number of
degrees of freedom. What is affordable is to \emph{probe} it, by evaluating the
leading omitted bands --- one shared stack serving every candidate degree ---
and to complete the amplitude from the model's measured degree-variance
spectrum.

\paragraph{The plan and the probe have different lifetimes.} The cancellation
target is an accumulated quantity and is stable along an arc; the local omitted
vector is field texture and is not. The same resolution scale that rules out a
tabulated direction also says how far ahead one may look: over a single
decision interval the vehicle's own state prediction is far more accurate than
$\pi r/N$, so probing forward along it is admissible, while over a seven-day
pilot arc it is not. A controller that plans offline on a cheap pilot arc and
probes forward one interval at a time respects that split, reads neither the
reference arc nor the field along it, and is the object the
campaign of Section~\ref{sec:campaign} evaluates.

Throughout, the budget is a common \emph{ceiling} on realized total quadratic
work, with the controller's own overhead charged against it, and not a nominal
per-call count. A ceiling rather than an equality because the objective is not
monotone in degree (Section~\ref{sec:algorithm}), so realized work is reported
beside every error rather than assumed identical. The
previous paper measured that the two differ by a median \ph{pct} and by a
factor of \ph{value} in the compute-starved regime, which is precisely the
regime where a budget matters, so the weaker definition is not used here to
carry a verdict.

\dnote{If the campaign lands in the amber band --- benchmark gap real,
controller cannot capture it --- this introduction inverts: the three
contributions become (i) the non-separable formulation and its certificate,
(ii) the resolution-scale argument as a measured \emph{cost} statement about
surrogates for the omitted direction, and (iii) a measured account of what deployable rules leave
on the table. The formulation and the decorrelation argument survive either
way; only the third paragraph changes.}

\section{Related work}
\label{sec:related}

\subsection{Truncation-degree selection}

\ph{one paragraph} Attenuation-only thresholds, degree-variance criteria in the
Meissl / Lowes--Mauersberger form \citep{RummelVanGelderen1995,Lowes1966}, and
Kaula-type spectral rules \citep{kaula1966theory}. Radius-driven selection
during propagation \citep{atallah2022radial}. These methods commonly select on
a local criterion --- force error, radius, spectral tail, or a state or
uncertainty proxy --- evaluated pointwise. The distinction drawn here is not
that such criteria are unknown to the literature but that none of them
allocates the \emph{degree history} jointly over a finite horizon under one
budget while retaining the signed cross-epoch coupling of
Eq.~\eqref{eq:objective}; that is the gap Section~\ref{sec:problem}
formalizes.

\subsection{Budgeted allocation}
\label{sec:related-allocation}

The separable case is classical. Allocating a scalar budget over independent
items through a single multiplier is Everett's generalized Lagrange multiplier
method \citep{everett1963generalized}, and the discrete form used for a set of
quantizers is the standard treatment in rate--distortion-style bit allocation
\citep{shoham1988efficient}, which also records that over a discrete set the
multiplier sweep recovers only supported efficient solutions, so a duality gap
may remain. The previous paper used exactly that construction at force level.

What changes here is that the objective couples the items, so neither result
applies directly. The coupled form of Eq.~\eqref{eq:allocation} is an integer
quadratic program with a single knapsack constraint, a class with its own
literature \ph{position against the quadratic-knapsack and integer-QP
literature; the citation pass supplies the references}. What that literature
supplies is the expectation of what is achievable --- a good feasible point and
a relaxation bound rather than a certified optimum --- and that is what
Section~\ref{sec:algorithm} delivers. What it does not supply is the
$\max(i,k)$ structure of Eq.~\eqref{eq:Q-structure}, which is specific to the
objective being a quadrature over prefix sums and is what makes a sweep linear
rather than quadratic in the number of epochs. From the separable literature,
what carries over is the language: an \emph{attainable} allocation reported as
such rather than as a bound.

\subsection{Goal-oriented error control}

Weighting a local residual by the sensitivity of a functional of interest is
the dual-weighted-residual idea in adaptive finite elements
\citep{becker2001optimal}, and the distinction it draws --- between a residual
norm and an adjoint-weighted residual --- is the same distinction the previous
paper measured in this setting. Two things separate the present problem from
that literature. The functional here is quadratic in an integral of the
residual rather than linear in it, so the weighting is not a scalar field;
and the control variable is a discrete degree with a quadratic cost, so the
adaptivity decision is a budgeted integer allocation rather than a refinement
flag.

\subsection{Discrete optimal control}

The structure this paper exploits is, read from the continuous side, an
adjoint one, and it is worth saying so rather than leaving a reader to notice
it. Treating the degree as a discrete control and the accumulated displacement
as a state, the allocation is a Mayer--Bolza problem with a mixed integer
control; Pontryagin's principle \citep{pontryagin1962mathematical} then
reduces the choice at each instant to minimizing a Hamiltonian against a
costate, and the cross term $\mathbf c_i$ of Eq.~\eqref{eq:cross-term} is
proportional to that costate transported to $t_0$ --- one half of
$\nabla_{\mathbf u_i}J=2\mathbf c_i$ under the normalization used here, the
factor depending on the adjoint convention. The forward--backward iteration of
Section~\ref{sec:controller-ifbda} is the corresponding two-point boundary
value problem solved by successive sweeps
\citep[Ch.~2]{bryson1975applied}, and the same reading places the method in
the mixed-integer optimal control literature, where relaxation followed by
sum-up rounding is the standard device \citep{sager2012integer}.

Naming the connection costs nothing and clarifies what is and is not new here.
The adjoint is not: it is a century old. What is specific is the combination
--- a spherical-harmonic truncation degree as the control, an objective that
is quadratic in a \emph{signed} integral so that the coupling is genuinely
cross-epoch, a discretization whose coupling matrix collapses onto
$\max(i,k)$ and therefore admits a certified bound rather than only a
stationarity condition, and a deployable approximation whose limiting
difficulty is informational rather than computational.

\subsection{Sensitivity in astrodynamics}

The state-transition matrix and the variational equations that generate it are
standard \citep[Ch.~4]{tapley2004statistical}. Mapping geopotential
coefficients to radial, transverse and normal position error analytically about
a Kepler orbit is classical \citep{rosborough1987radial}. \ph{one or two
sentences} on what is different here: the mapping is used forward, as an
allocation weight, and it is driven by a truncation defect rather than by an
estimation residual or a single coefficient.

\subsection{Multifidelity models}

\ph{short paragraph} \citep{peherstorfer2018multifidelity,
jones2019multifidelity}. The degree schedule is a fidelity schedule, and the
question of where to spend fidelity under a budget is shared; what is specific
here is that the low-fidelity model's error has an exploitable sign structure
along the trajectory.

\section{The allocation problem}
\label{sec:problem}

\subsection{Objective}

Let $\mathbf x_{\mathrm{ref}}(t)$ be a reference trajectory on $[0,T]$ and let

\begin{equation}
  \Delta\mathbf a(t,N) \;=\;
  \mathbf a_N\bigl(\mathbf x_{\mathrm{ref}}(t)\bigr)
  -\mathbf a_{\mathrm{ref}}\bigl(\mathbf x_{\mathrm{ref}}(t)\bigr)
  \label{eq:defect}
\end{equation}

be the truncation acceleration defect of degree $N$ measured along it. To first
order the induced \emph{state} perturbation is
$\delta\mathbf x(t)=\int_0^t \Phi(t,\tau)\,\mathbf B\,\Delta\mathbf a(\tau)\,\dd\tau
\in\mathbb{R}^6$, with $\Phi$ the reference state-transition matrix and
$\mathbf B$ mapping acceleration into the velocity block
\citep[Ch.~4]{tapley2004statistical}. The metric this paper reports is a
\emph{position} RMS, so the position block must be selected before the norm is
taken. With $\mathbf H_r=[\,\mathbf I_3\;\;\mathbf 0\,]$,

\begin{equation}
  E^2 \;=\; \frac{1}{T}\int_0^{T}
  \Bigl\lVert \mathbf H_r\!\int_0^{t}\Phi(t,\tau)\,\mathbf B\,
  \Delta\mathbf a(\tau)\,\dd\tau \Bigr\rVert^2 \dd t .
  \label{eq:objective}
\end{equation}

Taking the norm of the full six-vector would mix position and velocity units
and would not be the quantity the campaign measures; $\mathbf H_r$ is carried
explicitly through every expression below for that reason.

The decision variable is the degree history $N(\cdot)$, and the cost is the
quadratic work $\int N(t)^2\,\dd t$ it implies.

\subsection{Why the separable treatments do not apply}

Equation~\eqref{eq:objective} takes the norm \emph{after} the integral, and
that is the whole of the difficulty. Once the objective is discretized, Section~\ref{sec:problem-discretization}
puts it in the form $\mathbf u^{\top}\mathbf Q\,\mathbf u$, and the three
functionals separating this paper from its predecessor are then literally three
readings of \emph{one} matrix:

\begin{equation}
  \underbrace{\textstyle\sum_i \lVert\Delta\mathbf a_i\rVert^2\Delta t_i}
             _{\text{(i) defect norm: no }\mathbf Q}
  \qquad
  \underbrace{\textstyle\sum_i \Delta t_i\,
     \Delta\mathbf a_i^{\top}\mathbf K_i\,\Delta\mathbf a_i}
             _{\text{(ii) diagonal kernel}}
  \qquad
  \underbrace{\mathbf u^{\top}\mathbf Q\,\mathbf u}
             _{\text{(iii) all of }\mathbf Q}
  \label{eq:three-functionals}
\end{equation}

Functionals (i) and (ii) are separable in the per-epoch degrees --- (ii)
because a block-diagonal quadratic form decouples --- so a single multiplier
decides each epoch independently \citep{everett1963generalized}. Functional
(iii) is not: changing $N_i$ changes the sum against which every other epoch's
contribution is measured.

Writing the ladder this way removes an arbitrary choice: the sensitivity weight
in (ii) is not borrowed from anywhere, it is the objective's own coupling
evaluated on its diagonal. \emph{RQ2 is therefore exactly: how much of the gain
lives off the diagonal?} The previous paper's finding --- that a signed,
remaining-horizon-weighted statistic still ranks two policies the wrong way
round on \ph{n} of \ph{n} orbits --- says the answer is not zero.

\paragraph{The diagonal restriction has to be taken in the kernel, not in the
matrix.} Written continuously, $J$ is a double integral against a kernel,
$J=\int\!\!\int\Delta\mathbf a(\tau)^{\top}\mathbf K(\tau,\sigma)
\Delta\mathbf a(\sigma)\,\dd\tau\,\dd\sigma$, and the set $\tau=\sigma$ has
measure zero. ``Dropping the off-diagonal blocks'' of the discrete
$\mathbf Q$ is therefore not a definition one can take literally: since
$\mathbf u_i\propto\Delta t_i$, the raw diagonal term
$\mathbf u_i^{\top}\mathbf Q_{ii}\mathbf u_i$ carries $\Delta t_i^2$ and shrinks
as the grid is refined, so a criterion built on it would depend on the
accumulation resolution --- which, with the refined grid of
Section~\ref{sec:problem-grids}, is not even uniform along the arc.

What is well defined is the kernel \emph{on} the diagonal, a local sensitivity
$\mathbf K(t,t)$:

\begin{equation}
  \mathbf K(t,t) \;=\; \frac{1}{T}\int_t^T
  \mathbf B^{\top}\Phi(\tau,t)^{\top}\mathbf H_r^{\top}\mathbf H_r\,
  \Phi(\tau,t)\,\mathbf B\;\dd\tau ,
  \qquad
  \mathbf K_i \;=\; \frac{1}{T}\,\mathbf B^{\top}\Phi(t_0,m_i)^{\top}
  \mathbf A_i\,\Phi(t_0,m_i)\,\mathbf B ,
  \label{eq:local-sensitivity}
\end{equation}

the discrete form following from Eq.~\eqref{eq:Q-structure} and the cocycle
property. Functional (ii) is then $\sum_i\Delta t_i\,
\Delta\mathbf a_i^{\top}\mathbf K_i\Delta\mathbf a_i$, one power of $\Delta t$
per term like (i), and it converges to a physical limit under refinement.
Equivalently it is $\sum_i \mathbf u_i^{\top}\mathbf Q_{ii}\mathbf u_i/\Delta t_i$
--- the raw diagonal with the extra $\Delta t$ removed.

Nothing else moves. $\mathbf K_i$ is built from the same suffix sequence
$\mathbf A_i$, so it costs nothing extra; and (i) and (iii) were already
grid-consistent, (iii) because a double sum over $\mathbf u$ carries
$\Delta t_i\Delta t_k$ and is a genuine double integral. The derivation of
$\mathbf Q$, the prefix--suffix structure and the Frank--Wolfe certificate are
untouched.

\subsection{Discretization}
\label{sec:problem-discretization}

Fix an accumulation grid of $M$ cells with \emph{edges}
$e_0=0<e_1<\dots<e_M=T$ and widths $\Delta t_i=e_{i+1}-e_i$ --- not assumed
uniform, for the reason given in Section~\ref{sec:problem-grids}. Two index
sets are in play and they are deliberately distinct: $i=0,\dots,M-1$ runs over
\emph{cells}, whose midpoints $m_i$ carry the defect, and $j=0,\dots,M$ runs
over \emph{edges}, where the objective is sampled. Keeping them apart is what
the next paragraph is about; conflating them puts a term in the coupling that
violates causality.

With $\mathbf H_r=[\,\mathbf I_3\;\;\mathbf 0\,]$ selecting the position block,

\begin{equation}
  \mathbf u_i(N) = \Phi(t_0,m_i)\,\mathbf B\,\Delta\mathbf a(m_i,N)\,
  \Delta t_i ,
  \qquad
  \mathbf S_j = \sum_{i\le j-1}\mathbf u_i,
  \qquad
  \delta\mathbf r(e_j) = \mathbf H_r\,\Phi(e_j,t_0)\,\mathbf S_j ,
  \label{eq:u-and-S}
\end{equation}

using only the cocycle property $\Phi(e_j,m_i)=\Phi(e_j,t_0)\Phi(t_0,m_i)$.
The prefix runs to $j-1$ because cell $i$ occupies $[e_i,e_{i+1}]$ and is
therefore felt at edge $e_j$ exactly when $j\ge i+1$: at its own left edge the
impulse has not yet happened, and $\mathbf S_0=\mathbf 0$.

\paragraph{Which quadrature rule the prefix sum permits.} That
$\mathbf S_j$ is a \emph{plain} prefix sum is not a notational convenience; it
requires the inner weight of epoch $i$ to be independent of the outer epoch
$j$, and that requirement selects the rule. A trapezoid rule for the inner
integral has endpoint coefficient $(t_j-t_{j-1})/2$, which depends on $j$, so
carrying it would replace Eq.~\eqref{eq:Q-structure} by that structure plus a
$j$-dependent correction and forfeit the $\max(i,k)$ form. A rectangle rule
keeps the structure but is first-order, which at two samples per correlation
time is not a rounding detail.

We therefore place the accumulation nodes at the \emph{midpoints} of the grid
cells and sample the outer integral at the cell \emph{edges}. The inner weight
is then the full cell width --- independent of $j$, so the $\max(i,k)$
structure survives --- while the rule itself is the midpoint rule and
second-order, and $\mathbf S_j$ is the exact prefix sum of that
discretization through edge $j$. The outer integral carries no prefix sum and
is free to use the trapezoid rule on the edges. Both weight sets are
non-negative, which is what Section~\ref{sec:problem-structure} needs.

The price of separating the two grids is that the suffix in
Eq.~\eqref{eq:Q-structure} starts one index later than it would on a single
grid, and that offset is not cosmetic: it is the difference between a
discretization in which a cell influences the state \emph{after} it and one in
which it also influences the state at its own left edge.

The outer integral of Eq.~\eqref{eq:objective} is likewise a quadrature, with
weights $\omega_j\ge0$ summing to $T$, so the discretized objective is a
quadratic form in the stacked vector
$\mathbf u=(\mathbf u_0,\dots,\mathbf u_{M-1})$,

\begin{equation}
  J(\mathbf u) \;=\; \frac{1}{T}\sum_{j} \omega_j\,
  \mathbf S_j^{\top}\mathbf M_j\mathbf S_j
  \;=\; \mathbf u^{\top}\mathbf Q\,\mathbf u ,
  \qquad
  \mathbf M_j \;=\; \Phi(t_j,t_0)^{\top}\mathbf H_r^{\top}\mathbf H_r\,
  \Phi(t_j,t_0) ,
  \label{eq:quadratic-form}
\end{equation}

so that each $\mathbf M_j$ is a symmetric positive-semidefinite $6\times6$
block and no dimensional ambiguity is left in the notation. On a uniform grid
$\omega_j=\Delta t$ and $T=M\Delta t$, and Eq.~\eqref{eq:quadratic-form} reduces to
the unweighted average $M^{-1}\sum_j\mathbf S_j^{\top}\mathbf M_j\mathbf S_j$;
every expression below is written in the weighted form so that the two cases
are one derivation rather than two.

\paragraph{The decision variable lives on the decision grid.} The sum in
Eq.~\eqref{eq:quadratic-form} runs over accumulation epochs, but the degree
cannot be chosen $M$ times. Let $g(i)=q$ map an accumulation epoch to its
decision interval, $I_q=\{i: g(i)=q\}$, and write

\begin{equation}
  N_i = N_{g(i)} ,
  \qquad
  W_q \;=\; \sum_{i\in I_q}\Delta t_i
  \label{eq:interval-weight}
\end{equation}

for the interval's total time weight. The allocation problem is then

\begin{equation}
  \min_{(N_1,\dots,N_{K_{\mathrm{dec}}}) \,\in\, \mathcal{N}^{K_{\mathrm{dec}}}}
  \; \mathbf u^{\top}\mathbf Q\,\mathbf u
  \qquad \text{subject to} \qquad
  \sum_q W_q\,N_q^2 \;\le\; B \;=\; B_1 T ,
  \qquad K_{\mathrm{dec}} \ll M .
  \label{eq:allocation}
\end{equation}

The budget is a time integral of the instantaneous work rate, not a count of
samples: $B_1$ of Section~\ref{sec:setup-budgets} is the time average
$\langle N^2\rangle_t$, so the total over an arc of length $T$ is $B_1T$. Weighting by $W_q$ rather than by
the sample count $\lvert I_q\rvert$ is what keeps this correct when the
accumulation grid is refined near perilune --- an unweighted count would charge
the budget as though perilune occupied most of the arc simply because it is
sampled most densely. On a uniform grid $W_q=\lvert I_q\rvert\Delta t$ and the
two agree.

This is not only an implementability constraint. Optimizing over $M$ variables
would give the benchmark more \emph{switching freedom} than any deployable
controller has, and its advantage would then be partly a statement about grid
resolution rather than about trajectory-aware information. Every allocation
compared in this paper --- benchmark and controller alike --- is piecewise
constant on the same decision grid.

The result is an integer quadratic program with one knapsack constraint. Its
structure, however, is unusually favourable.

\subsection{The structure that makes it cheap}
\label{sec:problem-structure}

Because $\mathbf S_j$ is a prefix sum, the block of $\mathbf Q$ coupling cells
$i$ and $k$ collects exactly those edges that see both. Cell $i$ is seen at
edge $j$ when $j\ge i+1$, so the pair is seen when $j\ge\max(i,k)+1$:

\begin{equation}
  \mathbf Q_{ik} \;=\; \frac{1}{T}\!\!\sum_{j\ge\max(i,k)+1}\!\!\omega_j\,
  \mathbf M_j
  \;=\; \frac{1}{T}\,\mathbf A_{\max(i,k)} ,
  \qquad
  \mathbf A_i = \sum_{j= i+1}^{M} \omega_j\,\mathbf M_j .
  \label{eq:Q-structure}
\end{equation}

The offset by one is the whole content of separating cells from edges, and it
is worth stating why it cannot be dropped. Writing $\mathbf A_i=\sum_{j\ge i}$
instead would add $\omega_i\mathbf M_i$ to every suffix block, which credits
cell $i$ with displacing the state at $e_i$ --- before its own impulse. The
resulting form is not a small perturbation of the right one: on a random
five-cell instance it misstates $J$ by $26\%$, and it breaks the identity
between $\mathbf K_i$ and the diagonal of the kernel.

The quadrature weights enter only inside the suffix sequence, so the structure
that makes the problem tractable is untouched by a non-uniform grid. They are
required to be non-negative, which the trapezoid rule used here satisfies:
with $\omega_j\ge0$ and each
$\mathbf M_j$ positive semidefinite, $\mathbf Q$ is positive semidefinite, and
that is what the convex relaxation of Section~\ref{sec:algorithm-certificate}
rests on.

$\mathbf Q$ therefore has no free entries at all: it is generated by the single
suffix sequence $\mathbf A_i$, whose members are symmetric positive-semidefinite
$6\times6$ blocks --- the same size as $\mathbf M_j$, since $\mathbf u_i$ is a
state perturbation and the position selection happens inside $\mathbf M_j$
rather than on $\mathbf u_i$. The cross term needed by any coordinate method
follows from one prefix sum and one suffix sum,

\begin{equation}
  \mathbf c_i \;:=\; \sum_{k}\mathbf Q_{ik}\mathbf u_k
  \;=\; \frac{1}{T}\Bigl[\;
    \mathbf A_i \sum_{k\le i}\mathbf u_k
    \;+\; \sum_{k>i}\mathbf A_k\,\mathbf u_k
  \;\Bigr] ,
  \label{eq:cross-term}
\end{equation}

so evaluating $\mathbf c_i$ at \emph{all} $M$ epochs costs $\mathcal{O}(M)$
rather than the $\mathcal{O}(M^2)$ a dense $\mathbf Q$ would require, and the
gradient $\nabla J = 2\mathbf Q\mathbf u$ is available at the same cost. A
coordinate \emph{sweep} additionally tries $\lvert\mathcal{N}\rvert$ candidates
per decision interval and so costs
$\mathcal{O}(M\lvert\mathcal{N}\rvert)$; the two figures are quoted separately
throughout because they measure different things.

Grouping the epochs into decision intervals does not spoil this. For
$I_q=\{i_1<\dots<i_m\}$ the within-block term telescopes, because
$\max(i_p,i_{p'})=i_{p}$ whenever $p'<p$:

\begin{equation}
  \sum_{i,k\in I_q}\mathbf u_i^{\top}\mathbf Q_{ik}\mathbf u_k
  \;=\;
  \frac{1}{T}\Biggl[
  \sum_{p}\mathbf u_{i_p}^{\top}\mathbf A_{i_p}\mathbf u_{i_p}
  \;+\; 2\sum_{p}\Bigl(\sum_{p'<p}\mathbf u_{i_{p'}}\Bigr)^{\!\top}
  \mathbf A_{i_p}\mathbf u_{i_p}\Biggr] ,
  \label{eq:within-block}
\end{equation}

which one running prefix sum evaluates in $\mathcal{O}(m)$, and the
between-block term is Eq.~\eqref{eq:cross-term} with the in-block part
subtracted. A full sweep over all $K_{\mathrm{dec}}$ decision variables and all
candidate degrees therefore remains
$\mathcal{O}(M\lvert\mathcal{N}\rvert)$. Section~\ref{sec:algorithm} uses this.

The terminal-position objective is the special case in which the outer
quadrature collapses onto the final edge, $\omega_j = T\,\delta_{jM}$, so that
$\mathbf A_i = T\,\mathbf M_M$ for every cell and

\begin{equation}
  J_T \;=\; \bigl\lVert\,\mathbf H_r\,\Phi(e_M,t_0)\textstyle\sum_i\mathbf u_i
  \,\bigr\rVert^2 ,
  \label{eq:terminal}
\end{equation}

with the transport matrix retained because the $\mathbf u_i$ have been carried
back to $t_0$; dropping it would score the accumulated perturbation in the
wrong frame. Everything below covers this case unchanged.

\subsection{Two grids, not one}
\label{sec:problem-grids}

The decision variable and the integrand live on different time scales, and
conflating them is a measurement error rather than a modelling choice.

The degree schedule $N(t)$ varies on the orbital time scale. The defect
$\Delta\mathbf a$ does not: its direction is set by the leading omitted
harmonics, whose half-wavelength at the spacecraft's radius is $\pi r/N$, so
along track it decorrelates over
\begin{equation}
  \tau_{\mathrm{corr}} \;\approx\; \frac{\pi r}{N\,v} .
  \label{eq:decorrelation}
\end{equation}
A signed integral sampled coarser than that is aliased, even where the
corresponding magnitude statistic is converged --- and the previous paper
verified convergence only of the magnitude, which moved by \ph{pct} between a
120~s and a 10~s cadence.

The speed in Eq.~\eqref{eq:decorrelation} is the inertial $v$, and that is a
deliberately conservative proxy. What carries the vehicle across the field's
texture is the body-fixed transverse rate $r\lVert\dot{\hat{\mathbf r}}_{\mathrm{BF}}\rVert$,
which equals $v$ only where the motion is transverse; on an eccentric arc the
radial component inflates $v$ without moving the ground track, so
Eq.~\eqref{eq:decorrelation} runs short and the grid and probe are costed
high rather than low. The convergence sweep of Section~\ref{sec:campaign}
settles the question empirically --- it varies the sampling and asks where the
schedule stops moving --- so nothing rests on the proxy being tight.

Equation~\eqref{eq:decorrelation} is not one number. Both $r$ and $v$ swing by
factors of several on the eccentric orbits studied here, and $N$ with them, so
$\tau_{\mathrm{corr}}$ runs from \ph{value}~s near a 50~km perilune at $N=300$
to \ph{value}~s near apolune, a range of about \ph{value}. A uniform
accumulation grid fine enough for perilune is therefore mostly spent where
nothing happens, and one coarse enough for apolune aliases the passage that
matters. The grid is refined on $\tau_{\mathrm{corr}}$ instead, which also
requires the budget to be accumulated with time weights,
$B_1 = \sum_i N_{g(i)}^2\,\Delta t_i / T$; on a uniform grid this reduces to
the mean of Section~\ref{sec:setup-budgets}. \wpref{WP1, WP4}

We therefore separate

\begin{itemize}
  \item the \emph{accumulation grid} $\Delta t_{\mathrm{acc}}$, on which
    $\Delta\mathbf a$, $\Phi$ and the partial sums are formed, chosen fine
    enough that $J$ and the recovered schedule are converged
    (Section~\ref{sec:setup-grids}); and
  \item the \emph{decision grid} $\Delta t_{\mathrm{dec}}$, on which $N$ is held
    constant, chosen coarse enough to be implementable.
\end{itemize}

A schedule is thus piecewise constant on the decision grid while its
consequences are integrated on the accumulation grid. Section~\ref{sec:setup-grids}
reports the convergence of both, and Section~\ref{sec:ablations} reports what
coarsening the decision grid costs. \wpref{WP4}

\section{Solving the allocation, and certifying it}
\label{sec:algorithm}

Problem~\eqref{eq:allocation} is solved twice: once from above, by a descent
that returns an attainable schedule, and once from below, by a relaxation that
bounds what any schedule could achieve. The pair is what licenses the language
used for the result.

\subsection{Attainable schedule: budgeted coordinate descent}
\label{sec:algorithm-descent}

Introduce the multiplier $\lambda$ on the budget and minimize
$J(\mathbf u)+\lambda\sum_q W_q N_q^2$ by sweeping the \emph{decision
intervals}. Holding all other intervals fixed, the terms involving $N_q$ are

\begin{equation}
  N_q \;\leftarrow\; \argmin_{N\in\mathcal{N}}\;
  \Bigl[\;
    \underbrace{\sum_{i,k\in I_q}\mathbf u_i(N)^{\top}\mathbf Q_{ik}\,
      \mathbf u_k(N)}_{\text{within block, Eq.~\eqref{eq:within-block}}}
    \;+\; 2\sum_{i\in I_q}\mathbf u_i(N)^{\top}\mathbf c_i^{(-q)}
    \;+\; \lambda\,W_q\,N^2
  \;\Bigr] ,
  \label{eq:coordinate-update}
\end{equation}

with the interval time weight $W_q$ of Eq.~\eqref{eq:interval-weight} in place
of the sample count, and
$\mathbf c_i^{(-q)}=\mathbf c_i-\sum_{k\in I_q}\mathbf Q_{ik}\mathbf u_k$
from Eq.~\eqref{eq:cross-term}. The first group of terms is the block's own
contribution --- a magnitude, which is what a separable method would see ---
and the second is the coupling: it rewards a defect pointing \emph{against} the
accumulated displacement. That second term is the whole difference between this
and a sensitivity-weighted allocation.

\begin{algorithm}[!htbp]
\caption{Budgeted coordinate descent for Eq.~\eqref{eq:allocation}}
\label{alg:descent}
\begin{algorithmic}[1]
\Require tabulated $\mathbf u_i(N)$ for $i=1..M$, $N\in\mathcal{N}$; suffix
  blocks $\mathbf A_i$; grouping $g$; budget $B$; start set $\mathcal{S}$
  (Table~\ref{tab:starts})
\State $\lambda \gets$ bracket from the separable solution
\Repeat \Comment{outer: bisection on the budget}
  \For{each start $s\in\mathcal{S}$} \Comment{multi-start, \emph{not} one start}
    \State $N \gets s$
    \Repeat \Comment{inner: sweeps}
      \State refresh prefix sums of $\mathbf u$ and suffix sums of
        $\mathbf A_k\mathbf u_k$
      \For{$q = 1 \dots K_{\mathrm{dec}}$}
        \State update $N_q$ by Eq.~\eqref{eq:coordinate-update}; patch the two
          running sums over $I_q$
      \EndFor
    \Until{no interval changes, or sweep limit}
    \State record $(J(N),\,N)$
  \EndFor
  \State $N^\star(\lambda) \gets \argmin_s J$ \Comment{best of starts, at this
    $\lambda$}
  \State $\lambda \gets$ bisect on $\sum_q W_q N_q^{\star2} - B$
\Until{budget met to tolerance \ph{value}}
\State \Return $N^\star$, $J(N^\star)$, and the spread of $J$ across
  $\mathcal{S}$
\end{algorithmic}
\end{algorithm}

The multi-start loop sits \emph{inside} the bisection rather than outside it:
each $\lambda$ is solved from scratch from all starts. That is more expensive
than warm-starting from the previous $\lambda$, and it is deliberate --- a
warm start would make the returned schedule depend on the bisection path, which
is precisely the failure mode the monotonicity check below is meant to detect.

Each sweep is $\mathcal{O}(M\lvert\mathcal{N}\rvert)$ by
Eqs.~\eqref{eq:cross-term} and~\eqref{eq:within-block} --- grouping changes the
number of decision variables but not the sweep cost --- and no field evaluation
occurs inside the loop: the $\mathbf u_i(N)$ table is built once per orbit.

\paragraph{Summation.} The prefix sums are accumulated in compensated
(Neumaier) arithmetic. This is not routine caution: $\mathbf S_j$ is a sum of
order $10^5$ signed terms whose cancellation is the very quantity the method
exploits, so the summation is ill-conditioned by construction, and the ratio of
$\sum_i\lVert\mathbf u_i\rVert$ to $\lVert\mathbf S_M\rVert$ is reported per
orbit as the conditioning diagnostic. \wpref{WP1}

The objective is not convex in the integer variables, so the descent is
started from several points and the spread across starts is reported rather
than suppressed \wpref{WP3}.

\paragraph{The outer bisection needs a monotonicity check.} For a global
minimizer the attained work $W(\lambda)$ is non-increasing in $\lambda$, which
is what licenses bisecting on it. What Algorithm~\ref{alg:descent} returns is a
best-of-starts local solution, and nothing guarantees that the basin it lands
in varies monotonically. We therefore sweep $\lambda$ on a dense logarithmic
grid before any campaign run and verify that the best-of-starts $W(\lambda)$ is
non-increasing. Where it is, bisection is used. Where it is not, the bisection
is abandoned for that orbit and the schedule is chosen as the feasible grid
point with the smallest $J$, which is a weaker procedure but not an unsound
one; the orbits on which this happens are counted and reported. \wpref{WP3}

\begin{table}[!htbp]
  \centering\small
  \caption{Starting schedules for Algorithm~\ref{alg:descent} and the spread of
  the attained objective across them. Declared before the runs.}
  \label{tab:starts}
  \begin{tabular}{@{}llc@{}}
    \toprule
    Start & Rationale & Median $J/J_{\mathrm{best}}$ \\
    \midrule
    \Fop    & constant degree at the budget & \ph{value} \\
    \Rrad   & budget-calibrated radial endpoint & \ph{value} \\
    \Rint   & interior member & \ph{value} \\
    \Asens  & separable sensitivity-weighted solution & \ph{value} \\
    random $\times\,8$ & spread & \ph{value}--\ph{value} \\
    \bottomrule
  \end{tabular}
\end{table}

\subsection{Certificate: convex-hull relaxation and a Frank--Wolfe gap}
\label{sec:algorithm-certificate}

Calling the descent result a benchmark requires knowing how far it can be from
the best possible schedule. Relax each \emph{decision interval} to the convex
hull of its candidate contributions,

\begin{equation}
  \mathbf u_i \;=\; \sum_{N\in\mathcal{N}}\theta_{g(i)N}\,\mathbf u_i(N),
  \qquad
  \theta_{qN}\ge0, \quad \sum_N \theta_{qN}=1,
  \qquad
  \sum_q W_q \sum_N \theta_{qN}N^2 \;\le\; B .
  \label{eq:relaxation}
\end{equation}

Every integer schedule is feasible here with $\theta$ an indicator and attains
the same objective, so the relaxed optimum $J^{*}_{\mathrm{rel}}$ is a lower
bound on the integer optimum $J^{*}$. The relaxed problem minimizes the convex
quadratic $J$ over a convex set, and its linear subproblem

\begin{equation}
  \mathbf s \;\in\; \argmin_{\theta \text{ feasible}}\;
  \langle \nabla J(\mathbf u), \mathbf u(\theta)\rangle
  \label{eq:fw-subproblem}
\end{equation}

decouples across decision intervals under a single multiplier on the budget. It
therefore has the same epoch-separable, single-multiplier \emph{structure} as
the sensitivity-weighted allocation of Section~\ref{sec:benchmark-sens}, but it
is not the same problem: its per-interval objective is
$\langle\nabla_q J,\mathbf u_q(N)\rangle$ rather than
$\sum_{i\in I_q}\Delta t_i\Delta\mathbf a_i^{\top}\mathbf K_i\Delta\mathbf a_i$,
and
$\nabla J = 2\mathbf Q\mathbf u$ depends on the current relaxed iterate.
Frank--Wolfe therefore applies with an
$\mathcal{O}(M\lvert\mathcal{N}\rvert)$ iteration
\citep{frank1956algorithm,jaggi2013revisiting}, and its duality gap supplies
the bound directly:

\begin{equation}
  J(\mathbf u) + \langle\nabla J(\mathbf u),\,\mathbf s-\mathbf u\rangle
  \;\le\; J^{*}_{\mathrm{rel}} \;\le\; J^{*}
  \;\le\; J(\mathbf u_{\mathrm{descent}}) .
  \label{eq:certificate}
\end{equation}

Write $L_{\mathrm{FW}}$ for the best (largest) value the left-hand side attains
over the iteration budget and $J_{\mathrm{desc}}$ for the descent objective.
Because the objective is a squared error, a gap in $J$ and a gap in $E$ are not
the same number, so both are defined and the threshold is fixed on the second:

\begin{equation}
  g_J \;=\; \frac{J_{\mathrm{desc}}-L_{\mathrm{FW}}}{J_{\mathrm{desc}}},
  \qquad
  g_E \;=\; 1-\sqrt{\frac{L_{\mathrm{FW}}}{J_{\mathrm{desc}}}}
  \;=\; 1-\sqrt{1-g_J} .
  \label{eq:certified-gap}
\end{equation}

$g_E$ is the quantity the paper's claims are in --- it states that the
descent schedule's position error is within $g_E$ of the best error any
schedule of this class could achieve --- and the naming rule below is bound to
it. A $10\%$ error gap corresponds to $g_J\approx19\%$.

Two degenerate cases are handled by rule rather than by judgement. The
Frank--Wolfe bound is valid but can be vacuous early in the iteration, where
the left-hand side of Eq.~\eqref{eq:certificate} may fall below zero; since
$J\ge0$ always, $L_{\mathrm{FW}}$ is taken as
$\max\{0,\ \text{best bound}\}$. If $L_{\mathrm{FW}}=0$ when the iteration
budget is exhausted, the certificate is vacuous, no gap is reported for that
orbit, and the orbit is counted in a separate column rather than dropped.

\paragraph{The subproblem must be solved exactly.} The bound in
Eq.~\eqref{eq:certificate} is valid only if $\mathbf s$ attains the minimum of
Eq.~\eqref{eq:fw-subproblem} over the \emph{relaxed} feasible set. Stopping the
multiplier bisection at the nearest schedule whose integer budget falls below
$B$ does not do that: it returns a feasible point, not a minimizer, and the
resulting quantity is no longer a lower bound. Because
Eq.~\eqref{eq:relaxation} permits fractional mixing, the budget can be met with
equality --- at the critical multiplier the solution mixes the two adjacent
degrees of the affected intervals with the weights that saturate the
constraint. The implementation constructs that mixture explicitly and asserts
$\lvert\sum_q W_q\sum_N\theta_{qN}N^2 - B\rvert$ at machine
precision on every iteration; a run in which the assertion fails produces no
certificate.

\paragraph{What the certificate does and does not license.} It bounds the
distance to the best schedule \emph{of this problem}: the first-order
displacement along a fixed reference trajectory, on a fixed accumulation grid,
over a discrete degree set. It says nothing about the nonlinear trajectory,
which Section~\ref{sec:campaign-fixedpoint} tests separately. We adopt the
naming rule fixed before the runs: the quantity is called an \emph{oracle} only
if the median $g_E$ of Eq.~\eqref{eq:certified-gap} is below $0.10$ and no more
than \ph{n} orbits return a vacuous certificate; otherwise it is reported as a
\emph{linearized trajectory-aware allocation benchmark} and every ratio taken
against it is read as conservative, in the same sense the previous paper's
force-level Lagrangian benchmark was.

\dnote{One of the two names has to be chosen once Table~\ref{tab:certificate}
exists, and the choice is bound to the gap number rather than to the result.
If the gap is large, delete the word ``oracle'' from the whole manuscript ---
including the abstract --- rather than qualifying it in place.}

\begin{table}[!htbp]
  \centering\small
  \caption{Certified gap between the Frank--Wolfe lower bound and the
  coordinate-descent schedule, by design and budget. \ph{fill from WP7}}
  \label{tab:certificate}
  \begin{tabular}{@{}lccccccc@{}}
    \toprule
    Design & $\beta$ & $L_{\mathrm{FW}}$ & $J_{\mathrm{desc}}$ & $g_J$ & $g_E$
      & worst $g_E$ & vacuous \\
    \midrule
    A & $1.00$ & \ph{value} & \ph{value} & \ph{pct} & \ph{pct} & \ph{pct} & \ph{n} \\
    B & $1.00$ & \ph{value} & \ph{value} & \ph{pct} & \ph{pct} & \ph{pct} & \ph{n} \\
    \bottomrule
  \end{tabular}
\end{table}

\subsection{Cost of the solve}

\ph{one paragraph} Table build per orbit; sweep count to convergence;
Frank--Wolfe iterations to a useful gap; memory of the vector table at the
converged accumulation grid. None of this is charged to the flight budget ---
the benchmark is not flown --- but it decides what is feasible to run across a
campaign, and the controller of Section~\ref{sec:controller} inherits a
reduced form of it, which \emph{is} charged. \wpref{WP1, WP3, WP7}

\section{Experimental setup}
\label{sec:setup}

\subsection{Field, dynamics and reference}

\ph{one paragraph} GRAIL product and degree \citep{Konopliv2014GRAIL,
GrailGL1800FSHADR,Goossens2020}; synthesis kernel and normalization
\citep{holmes2002unified,pines1973uniform}; perturbation model and the
conventions it follows \citep{montenbruck2000satellite,vallado2013fundamentals};
third bodies and ephemeris \citep{park2021de440,acton2018spice}; integrator and
tolerances \citep{hairer1993ode,virtanen2020scipy}; lunar orientation
\citep{Mazarico2018}. Adopted \emph{reference} degree per orbit and the sense in
which every result below is model-relative: what is measured is what the adopted
expansion requires of its own truncations, not the distance to the true lunar
field.

\subsection{Populations}

\begin{table}[!htbp]
  \centering\small
  \caption{Orbit populations. Designs A and B are independent scrambles of the
  same scrambled-Sobol coverage construction \citep{joe2008sobol}; C was drawn
  and frozen after A and B were known to agree. That the design points are
  dependent by construction is why no significance level is attached to any
  tally. The wide-elliptic population lies outside the sampled factor box
  and brackets the elliptical subsatellite orbits flown on Kaguya
  \citep{kato2010kaguya}.}
  \label{tab:populations}
  \begin{tabular}{@{}llll@{}}
    \toprule
    Population & Orbits & Geometry & Role \\
    \midrule
    Design A & \ph{n} & \ph{box} & primary \\
    Design B & \ph{n} & \ph{box} & independent replication \\
    Design C & \ph{n} & \ph{box} & third design \\
    Wide-elliptic & \ph{n} & perilune \ph{km}, apolune \ph{km} & regime bound \\
    Strata (5) & \ph{n} each & inclination / argument of perilune restricted & geometry bound \\
    Low perilune & \ph{n} & \ph{km} & linearization bound \\
    \bottomrule
  \end{tabular}
\end{table}

\subsection{Budgets}
\label{sec:setup-budgets}

Four cost quantities appear and are never interchanged.

\begin{table}[!htbp]
  \centering\small
  \caption{Cost definitions. $B_2$ is the primary budget throughout; $B_1$ is
  reported only for comparability with the previous paper.}
  \label{tab:budget-definitions}
  \begin{tabular}{@{}llp{0.44\textwidth}@{}}
    \toprule
    Symbol & Quantity & Note \\
    \midrule
    $B_1$ & nominal work rate, the \emph{time average} $\langle N^2\rangle_t
            = \sum_i N_{g(i)}^2\Delta t_i/T$
          & available without propagating; does not survive contact with the
            integrator. On the uniform reference output grid this is the plain
            mean, which is how the previous paper reported it \\
    $B_2$ & realized total quadratic work $\sum_{\mathrm{RHS}}N^2$
          & \textbf{primary}; requires propagation \\
    $B_3$ & measured serial gravity-kernel time
          & architecture-dependent; idle machine, three repeats \\
    $B_+$ & controller overhead, in the same units as $B_2$
          & \textbf{included in every budget reported for \Cplan, \Clite\ and
            \Cfb}; conversion below \\
    \bottomrule
  \end{tabular}
\end{table}

\paragraph{What $B_+$ is measured in.} $B_2$ counts quadratic work,
$\sum_{\mathrm{RHS}}N^2$, and $B_+$ is added to it, so the two must carry the
same units. Two of the overhead terms already do: the pilot arc and the online
probe are gravity evaluations and enter directly as $\sum N^2$. The rest ---
the two-body predictor, the planning sweeps, the online $\argmin$ and its
bookkeeping --- are processor time and are converted at the measured rate of
the same run,

\begin{equation}
  B_{+,\mathrm{eq}} \;=\; t_+ \Bigl/ \frac{t_{\mathrm{kernel}}}{B_2} ,
  \label{eq:overhead-units}
\end{equation}

with $t_{\mathrm{kernel}}$ the arc's measured serial gravity-kernel time. The
ratio $t_{\mathrm{kernel}}/B_2$ is a seconds-per-unit-work coefficient, so
Eq.~\eqref{eq:overhead-units} expresses non-gravity time in $N^2$ units and
Eq.~\eqref{eq:matching} adds like to like. The coefficient is
architecture-dependent and is measured under the same conditions as $B_3$, on
an idle machine with three repeats; every table quoting $B_+$ names the machine
it was measured on.

\paragraph{Normalization.} $\beta$ is defined at the $B_1$ level only, as the
allocation \emph{target} $\beta=B_1/N_{\mathrm{crit}}^2$. The realized
counterpart is reported separately as
$\beta_2 = B_2/(K_{\mathrm{ref}}N_{\mathrm{crit}}^2)$, with
$K_{\mathrm{ref}}$ the comparator's right-hand-side call count on the same arc;
dividing $B_2$ by $N_{\mathrm{crit}}^2$ alone would fold the call count into
the same number and is not done.

\paragraph{Matching convention.} Comparisons are matched on a common work
\emph{ceiling}, and the controller's overhead comes out of its own allocation:

\begin{equation}
  B_2 + B_+ \;\le\; B_{\mathrm{tot}} \;=\; B_{2,\mathrm{F}} ,
  \label{eq:matching}
\end{equation}

with the inequality meant literally. Section~\ref{sec:algorithm-descent}
explains why an allocation can be better for spending less --- the objective
is not monotone in degree --- so the realized $B_2+B_+$ is reported beside
every error rather than assumed equal to the ceiling. A policy that wins under
a lower spend has won by more; one that wins only by saturating the ceiling is
reported as such.

with $B_+=0$ for every policy that is not a controller. No fictitious overhead
is added to the fixed comparator, so a controller that must plan and probe has
strictly less budget left for gravity evaluation than its comparator does ---
the conservative reading, and the one adopted throughout.

Equation~\eqref{eq:matching} is not complete until it says which side moves,
and the answer is not the one the previous campaign used. There, the comparator
was re-propagated to spend what the candidate had spent. That cannot be done
here, because the capture fraction of Section~\ref{sec:campaign} divides one
difference by another and both must be taken against \emph{the same}
$E_{\Fop}$; a comparator that slides to meet each candidate would put two
different $\Fop$ instances in the numerator and the denominator.

We therefore anchor. For each orbit and each $\beta$, $B_{\mathrm{tot}}$ is
declared in advance as the realized total quadratic work of $\Fop(\beta)$ on
that arc --- at $\beta=1$ that is the archived critical-degree run, so the
anchor already exists and is not recomputed --- and \emph{every candidate is
calibrated to it}.

That calibration costs propagations, and the cost is a consequence of the
previous paper's own result rather than an inefficiency: a schedule's realized
work is not known until it is flown, and calibrating on nominal per-call work
misses it by a median \ph{pct}. Each candidate is therefore propagated,
measured, its multiplier adjusted, and propagated again until it sits under
the ceiling without slack it does not want:
$B_2+B_+\le B_{\mathrm{tot}}$, and either
$(B_{\mathrm{tot}}-B_2-B_+)/B_{\mathrm{tot}}\le\ph{pct}$ or the schedule is
the unconstrained optimum of its own criterion. The second branch is what
keeps the calibration from forcing work onto a policy whose objective would
worsen for taking it. That takes \ph{value} propagations per candidate in
practice, and Section~\ref{sec:campaign} reports the achieved spend alongside
every verdict.

The propagated grid is $\beta\in\{0.50,0.75,1.00,1.25,1.50\}$;
$\beta=3$ is evaluated at benchmark level only, because at that budget the
previous campaign found \ph{n} of \ph{n} comparisons undecidable and the
experiment stops being an instrument.

\subsection{Grids and their convergence}
\label{sec:setup-grids}

\ph{one paragraph plus table} The campaign runs on an accumulation grid
refined on $\tau_{\mathrm{corr}}$, so the convergence parameter is not an
absolute step but the number of samples per correlation time,
$n_s=\tau_{\mathrm{corr}}(t)/\Delta t_i$, swept over
$\{0.5,1,2,4,8\}$. A uniform-grid sweep at
$\Delta t_{\mathrm{acc}}\in\{240,120,60,30,10\}$~s is run alongside it, both
because it is the comparable quantity in the previous paper and because it
isolates whether the refinement, rather than the resolution, is what matters.
Three quantities are checked in each: the objective, the recovered schedule
itself as a degree-profile distance, and verdict stability. The decision grid
is swept separately at $\Delta t_{\mathrm{dec}}\in\{60,120,300\}$~s. The
campaign runs at $n_s=\ph{value}$. \wpref{WP4}

\dnote{Section~\ref{sec:problem} predicts that the signed integral needs a
finer grid than the magnitude statistic did. If the panel contradicts that ---
if 120~s is already converged for $J$ and for the schedule --- say so plainly
and delete the prediction rather than keeping it as motivation.}

\subsection{Resolution rule and statistics}

A comparison is decided only when the two errors differ by more than their
summed reference-inclusive numerical envelope, $M_{\mathrm{res}}>1$; comparisons
below that are \emph{undecided} and are counted for neither side. No
significance level is attached to any tally: the points of a scrambled design
are dependent by construction and the resolution rule censors the sample
non-randomly, so a binomial test would have no valid null. Direction rests on
replication across independent populations. Ratios are medians of per-orbit
ratios, never ratios of medians.

\subsection{Comparators}
\label{sec:setup-comparators}

Two constant-degree comparators are carried, and they answer different
questions.

\begin{itemize}
  \item \Fop, the integer degree nearest the budget, is what a user can select
    in advance. Beating it is the claim.
  \item \Fenv, the lowest-error member of the constant family under the same
    budget, is found by measuring every candidate against the reference. It is
    a post-hoc lower envelope, not a policy; the previous paper measured it a
    median \ph{value} times better than the saturating degree. It is reported
    to bound how much of the constant family's potential \Fop\ leaves unused,
    and no hypothesis of this paper requires beating it.
\end{itemize}

\Rint, the interior member of the radial--constant interpolation family, is
carried as the strongest previously available deployable schedule.

\subsection{Pre-registration and its stopping rules}
\label{sec:setup-prereg}

Hypotheses, gates, sweeps, comparator rules, censoring and the naming rule of
Section~\ref{sec:algorithm-certificate} were written to
\ph{oa01/oa02/oa03 json} and hashed before any aggregate result was inspected.
Anything added later is labelled post hoc and reported separately; rejected
hypotheses are reported as rejected rather than reworded. \ph{outcome list}

Two of those rules bound the campaign rather than a single number, and both are
stated here because later sections invoke them.

\paragraph{Budget gate.} The first propagated campaign is at $\beta=1$ on one
design. The budget grid is extended to the rest of the populations only if the
controller also takes the resolved majority at $\beta=0.75$ \emph{and}
$\beta=1.25$ on a stratified sixteen-orbit subset. If it does not, the gain
belongs to one budget point, the population expansion is not run, and the
result is reported as such.

\paragraph{Screening escape.} The variational screen of
Section~\ref{sec:campaign} is calibrated on policies that were not built by
minimizing $J$, so its transfer to policies that were is an assumption rather
than a measurement. If the propagated sign disagrees with the screened sign,
the implementation is diagnosed; after two diagnostic rounds without
resolution the instrument is declared uncalibrated for this class, the screen's
selection is discarded, and the campaign proceeds on propagation alone with a
reduced candidate list. The cost of that path is reported rather than absorbed.
The rule exists so that a disagreement cannot become an open-ended loop.

\section{How much is attainable, and where it comes from}
\label{sec:benchmark}

This section answers RQ1 and RQ2. Nothing in it is propagated: every quantity
is computed from the tabulated defect and the reference state-transition
matrix, integrated with the variational equations along the reference arc,
so the comparison is free of integrator noise and of the step-size effects the
previous campaign had to control. What it therefore establishes is a
first-order statement about a fixed reference trajectory;
Section~\ref{sec:campaign} tests whether the propagated arcs agree.

\paragraph{Predicted error, not measured error.} Every error in this section is
the first-order prediction $\hat E=\sqrt{J}$ of Eq.~\eqref{eq:quadratic-form},
and every ratio is written $\hat\rho$. The propagated quantities $E$ and $\rho$
appear only from Section~\ref{sec:campaign} onward. The two are kept
typographically distinct throughout because they answer different questions and
are produced at different stages: no hypothesis stated on $\hat E$ is reported
as though it had been tested on $E$. Where the resolution rule is invoked here
it uses each orbit's archived numerical envelope, which is external to this
paper and pre-existing, since no propagation of this campaign has yet occurred
to supply one.

\subsection{The two separable allocations}
\label{sec:benchmark-sens}

Two of the four rungs below are separable and are stated here once, because
the Frank--Wolfe subproblem of Section~\ref{sec:algorithm-certificate} has the
same structure. Both are per-interval minimizations with a single multiplier
$\lambda$ bisected to the budget:

\begin{equation}
  \Aforce:\;\;
  \argmin_{N\in\mathcal{N}}\Bigl[\textstyle\sum_{i\in I_q}
    \lVert\Delta\mathbf a(t_i,N)\rVert^2\Delta t_i
    +\lambda W_q N^2\Bigr],
  \qquad
  \Asens:\;\;
  \argmin_{N\in\mathcal{N}}\Bigl[\textstyle\sum_{i\in I_q}
    \Delta t_i\,\Delta\mathbf a(t_i,N)^{\top}\mathbf K_i\,
    \Delta\mathbf a(t_i,N)
    +\lambda W_q N^2\Bigr].
  \label{eq:separable}
\end{equation}

\Aforce\ is the force-level benchmark of the previous paper: it never sees
$\Phi$ at all. \Asens\ is the natural sensitivity-weighted repair, and it is
worth being precise about where its weight comes from, because two wrong
answers are close at hand. Borrowing a scalar such as
$\lVert\Phi(T,t)\mathbf B\rVert^2$ from the terminal-state proxy would be an
arbitrary import, and matched to the wrong metric besides. Taking the discrete
diagonal $\mathbf u_i^{\top}\mathbf Q_{ii}\mathbf u_i$ literally would make the
criterion depend on the accumulation resolution, for the reason given in
Section~\ref{sec:problem-discretization}. The weight is instead the local
sensitivity kernel $\mathbf K_i$ of Eq.~\eqref{eq:local-sensitivity} --- the
objective's own coupling evaluated on its diagonal, grid-invariant, and matched
to the arc-RMS metric rather than to a terminal one.

The ladder is therefore a decomposition of one trajectory kernel and not a
comparison of three unrelated criteria: \Aforce\ ignores it, \Asens\ retains
its local diagonal $\mathbf K_i$, \Asign\ retains the full cross-epoch
coupling. RQ2 reduces to a single question --- how
much of the attainable gain lives off the diagonal --- and $G_{\mathrm{sens}}$
and $G_{\mathrm{sign}}$ below measure the two steps of that decomposition.

\subsection{The ladder}

Four allocations are compared at the same budget, each carrying one more piece
of the dynamics than the last.

\begin{table}[!htbp]
  \centering\small
  \caption{Allocation ladder at $\beta=1$. Entries are
  $\hat\rho(\Fop,\text{policy})=\hat E_{\Fop}/\hat E_{\text{policy}}$ on the
  first-order prediction, medians of per-orbit ratios, so values above one
  favour the policy. All four are evaluated on the same reference trajectories
  and the same tabulated defect, and none is propagated.
  \ph{fill from WP2, WP3}}
  \label{tab:ladder}
  \begin{tabular}{@{}llcccc@{}}
    \toprule
    & Allocation & Information used & A & B & C \\
    \midrule
    1 & \Rint\  & radius & \ph{value} & \ph{value} & \ph{value} \\
    2 & \Aforce & + reference field, pointwise & \ph{value} & \ph{value} & \ph{value} \\
    3 & \Asens  & + diagonal kernel $\mathbf K_i$ & \ph{value} & \ph{value} & \ph{value} \\
    4 & \Asign  & + signed coupling $\mathbf Q$ & \ph{value} & \ph{value} & \ph{value} \\
    \bottomrule
  \end{tabular}
\end{table}

\paragraph{RQ1 --- the attainable gain.} \ph{one paragraph} \Asign\ sits a
median \ph{value} above \Fop\ and \ph{value} above \Rint, the strongest
deployable schedule previously available. \ph{state whether the gain clears the
numerical envelope on how many orbits, and name the orbits where it does not.}

\dnote{Gate G1 of \code{../ROADMAP.md} is decided here. If
$\hat\rho(\Rint,\Asign)<1.5$ at the design median, this section becomes the
paper's principal result in the negative direction and
Sections~\ref{sec:controller} and~\ref{sec:campaign} contract sharply.}

\paragraph{RQ2 --- weight against sign.} The two increments are reported
separately because they are different claims. Raw error differences are not
comparable across orbits of different scale, so each increment is a log ratio:

\begin{equation}
  G_{\mathrm{sens}} = \log\frac{\hat E_{\Aforce}}{\hat E_{\Asens}},
  \qquad
  G_{\mathrm{sign}} = \log\frac{\hat E_{\Asens}}{\hat E_{\Asign}} ,
  \label{eq:increments}
\end{equation}

the first being what correcting the weighting buys and the second what the
coupling buys on top of it. \ph{one paragraph with the two medians and their
ratio.} The previous paper predicted that reweighting alone would not close the
gap; \ph{state whether that is confirmed, and by how much.}

\begin{figure}[!tbp]
  \centering
  \ph{figures/fig\_oa\_ladder.pdf}
  \caption{The allocation ladder, one point per orbit, designs pooled. Bars are
  medians. Rungs 1--2 use magnitudes only; rung 3 adds the sensitivity weight;
  rung 4 adds the signed coupling.}
  \label{fig:ladder}
\end{figure}

\subsection{What the schedules look like}

\ph{one paragraph plus figure} The signed allocation is not a monotone function
of radius, which is the point: at equal budget it \ph{describe the qualitative
structure --- where it spends relative to \Asens, and whether the departure is
phase-locked}. Reporting this matters because it is the part a deployable rule
has to imitate.

\begin{figure}[!tbp]
  \centering
  \ph{figures/fig\_oa\_schedules.pdf}
  \caption{Degree histories at $\beta=1$ on a representative orbit:
  \Fop, \Rint, \Asens\ and \Asign. The lower panel shows the accumulated
  displacement each produces.}
  \label{fig:schedules}
\end{figure}

\subsection{Certificate}

Table~\ref{tab:certificate} gives the certified gap. \ph{one paragraph}
\ph{State the median gap, the worst orbit, and therefore which of the two names
of Section~\ref{sec:algorithm-certificate} the paper uses.} Where the gap is
wide the reported ratios are conservative: the true optimum can only be better
than the schedule found, so the advantage over \Fop\ is understated and the
capture fraction of Section~\ref{sec:controller} overstated.

\subsection{Sensitivity of the benchmark to its own construction}

Three properties of the construction were audited and none moves a direction
reported above. \wpref{WP3, WP4}

\begin{itemize}
  \item \textbf{Start dependence.} Table~\ref{tab:starts}; the spread across
    \ph{n} starts is \ph{value}.
  \item \textbf{Accumulation grid.} Section~\ref{sec:setup-grids}; the schedule
    changes by \ph{value} between \ph{value}~s and \ph{value}~s.
  \item \textbf{Degree set.} \ph{one sentence} A coarser $\mathcal{N}$ can only
    make the benchmark worse, so the measured gap to the policies is
    conservative in that direction too.
\end{itemize}

\section{A deployable controller}
\label{sec:controller}

The benchmark of Section~\ref{sec:benchmark} evaluates the reference field at
every epoch of a reference trajectory it already knows. This section removes
that access and measures what each removal costs. This is RQ3.

\paragraph{What ``without the reference field'' means, precisely.} The
distinction is not between having and not having the gravity model --- a
propagator carries the coefficients by definition, and pretending otherwise
would be artificial. It is between \emph{global, offline model metadata} and
\emph{trajectory-specific evaluations of it}:

\begin{center}\small
\begin{tabular}{@{}lll@{}}
  \toprule
  & Available to the controller & Withheld \\
  \midrule
  Field & degree variances $P_n$ (a table of \ph{n} numbers) &
    $\Delta\mathbf a_{\mathrm{ref}}(t)$ along the arc \\
  Trajectory & a cheap pilot arc it computes itself & the reference arc \\
  Sensitivity & $\Phi$ from low-degree variational equations & the exact $\Phi$ \\
  \bottomrule
\end{tabular}
\end{center}

The controller may know the spectrum; it may not be told what the omitted
acceleration is at the epochs it will fly through. That is the information
boundary the rest of this section works against, and it is the one a real
mission faces.

\begin{table}[!htbp]
  \centering\small
  \caption{What the benchmark uses, and what replaces it.}
  \label{tab:substitutions}
  \begin{tabular}{@{}p{0.26\textwidth}p{0.34\textwidth}p{0.32\textwidth}@{}}
    \toprule
    Benchmark uses & Replaced by & Cost, charged to $B_+$ \\
    \midrule
    Reference field, pointwise & leading omitted bands (direction) plus the
      measured degree-variance spectrum (amplitude) & one synthesis per
      look-ahead probe point; measured \ph{pct} --- the dominant term \\
    Reference arc & cheap pilot arc, $N=40$, loose tolerance & measured
      \ph{pct} of the run's gravity work \\
    Exact $\Phi$ & low-degree variational equations integrated on the pilot
      arc, $\dot\Phi = \bigl[\begin{smallmatrix}\mathbf 0&\mathbf I\\
      \mathbf G_{40}&\mathbf 0\end{smallmatrix}\bigr]\Phi$ & \ph{pct} \\
    \bottomrule
  \end{tabular}
\end{table}

\subsection{The omitted acceleration has a direction, and tabulating it is not cheap}
\label{sec:controller-direction}

A tail criterion predicts $\lVert\Delta\mathbf a\rVert$. The coupling term of
Eq.~\eqref{eq:coordinate-update} needs $\Delta\mathbf a$ itself. A controller
supplied only with the magnitude can reach rung 3 of Table~\ref{tab:ladder} and
no further, by construction.

The direction varies on the model's own resolution scale. It is set by the
leading omitted harmonics, whose half-wavelength at the spacecraft's radius is
$\pi r/N$: about \ph{value}~km at $N=120$ and \ph{value}~km at $N=300$. A tabulation,
interpolation or fitted surrogate over $(r,\phi,\lambda)$ that does not itself
resolve that scale therefore cannot represent the direction, and one that does
resolve it carries a number of degrees of freedom comparable to the omitted
coefficients themselves, with the storage and evaluation cost that implies.
This is a statement about resolution and cost, not an impossibility claim about
any particular function class; Section~\ref{sec:ablations} measures it directly
with a surrogate built for the purpose and reports what it retains and what it
costs. \wpref{WP5}

\subsection{Band probe}

What is affordable is to evaluate the first few omitted bands and use their sum
as the direction:

\begin{equation}
  \hat{\mathbf v}(t,N)
  \;=\; -\,\gamma(N,r)\,\sum_{n=N+1}^{N+k}\mathbf a_n(t) ,
  \label{eq:probe}
\end{equation}

with $\gamma$ the ratio of the full omitted tail to the probed bands. Direction
comes from the probe; amplitude from $\gamma$.

$\gamma$ is taken from the model's \emph{measured degree-variance spectrum},
not from a fitted power law:

\begin{equation}
  \gamma(N,r) \;=\;
  \Biggl[\frac{\sum_{n>N}\sigma_a^2(n;r)}
              {\sum_{n=N+1}^{N+k}\sigma_a^2(n;r)}\Biggr]^{1/2} ,
  \label{eq:gamma}
\end{equation}

with $\sigma_a(n;r)$ the per-degree acceleration RMS built from the coefficient
degree variances $P_n$. The distinction matters because the lunar spectrum is
not a single power law --- the previous paper measured a spectral slope of
$2.13$ against an effective tail exponent of $1.76$ for the same field, and the
gap is a property of the spectrum's shape --- so a power-law extrapolation of
the tail carries a systematic error that a table of $P_n$ does not. And it
costs nothing in deployability: $P_n$ is roughly \ph{n} numbers, a
one-dimensional table.

That split is worth naming, because it is the paper's information statement in
its sharpest form. \textbf{What is cheap to carry is the spectrum; what is
expensive is the texture.} The amplitude of the omitted tail is a
one-dimensional function of degree and can simply be uploaded. Its direction is
a field on the sphere at resolution $\pi r/N$ and cannot; it has to be probed.

\paragraph{Cost, and why the cheap-looking figure is the wrong one.} At a point
where the field is \emph{already} being evaluated to degree $N$, extending to
$N{+}k$ costs $(N{+}k)^2-N^2\approx 2kN$ against $N^2$, that is $2k/N$ --- about
5\% at $N=120$, $k=3$ --- because the associated Legendre recursion is
incremental in degree and the probe inherits it.

That figure does not apply here, and saying so is the honest form of the cost
accounting. A look-ahead probe is taken at a point the integrator has \emph{not}
evaluated, and the recursion for $P_{nm}$ at a new location must run from
degree zero: obtaining band $n$ alone is not an $\mathcal{O}(n)$ operation in
any standard synthesis. \textbf{Each look-ahead probe point therefore costs
approximately one full synthesis}, not an increment on one.

Two properties of the construction do reduce the cost, but both act
\emph{within} a probe point rather than between them.

First, the band stack is \emph{shared} across candidates. Evaluating the individual bands
$\mathbf a_n$ once for $n$ up to $N_{\max}{+}k$ yields
$\hat{\mathbf v}(t,N)$ for \emph{every} candidate $N\le N_{\max}$ by partial
summation, at no further field evaluation. The candidates are therefore not
independent probes but partial sums of one stack.

Second, the candidate set is a window. The offline plan supplies a nominal
degree for each decision interval; writing $j(q)$ for its position in the
ordered candidate grid $\mathcal{N}=\{N_1<N_2<\dots\}$, the online decision
searches only

\begin{equation}
  \mathcal{N}_q \;=\; \{\,N_{j(q)-\delta},\,\dots,\,N_{j(q)+\delta}\,\} ,
  \label{eq:window}
\end{equation}

so that $\delta$ is a half-width in \emph{candidate-grid indices}, not in
degrees. Fixing it that way keeps the controller's complexity ---
$2\delta+1$ candidates per decision --- invariant when the grid is refined or
coarsened; the induced span in degrees follows from the grid and is reported
alongside the cost, since it is the span and not the index count that sets the
band stack. The stack runs from $N_{j-\delta}{+}1$ to $N_{j+\delta}{+}k$, so
the incremental cost is $\approx 2N(\Delta_{\mathrm{span}}+k)$ with
$\Delta_{\mathrm{span}}=N_{j+\delta}-N_{j-\delta}$, rather than $2kN$.

This is the plan/probe split made concrete: the plan sets the coarse
allocation, which needs the budget and the sensitivity and is smooth; the probe
corrects within a window, which needs the local texture. The window is a
declared parameter, swept in Section~\ref{sec:ablations}, and it handicaps the
controller relative to the benchmark --- which searches the full $\mathcal{N}$
--- so the cost of the restriction appears in the capture fraction rather than
being hidden.

\paragraph{What the overhead actually is.} Covering a decision interval takes
$n_{\mathrm{probe}}\approx\Delta t_{\mathrm{dec}}/\tau_{\mathrm{corr}}$ probe
points, and the interval contains $n_{\mathrm{RHS}}\approx
\Delta t_{\mathrm{dec}}\,r_{\mathrm{RHS}}$ right-hand-side calls, so the
overhead is

\begin{equation}
  \frac{n_{\mathrm{probe}}}{n_{\mathrm{RHS}}}
  \;\approx\; \frac{1}{\tau_{\mathrm{corr}}\;r_{\mathrm{RHS}}}
  \;=\; \frac{N v}{\pi r\, r_{\mathrm{RHS}}} ,
  \label{eq:probe-overhead}
\end{equation}

which is \emph{independent of the decision cadence}: probing twice as often
over intervals twice as short buys nothing.

Equation~\eqref{eq:probe-overhead} is a local rate, and the number that belongs
in a budget is its arc integral,

\begin{equation}
  \frac{B_{\mathrm{probe}}}{B_{\mathrm{tot}}}
  \;=\; \frac{1}{N_{\mathrm{RHS}}}\int_0^T \frac{\dd t}{\tau_{\mathrm{corr}}(t)}
  \;=\; \frac{1}{N_{\mathrm{RHS}}}\int_0^T
        \frac{N(t)\,v(t)}{\pi\,r(t)}\,\dd t ,
  \label{eq:probe-overhead-arc}
\end{equation}

and the two are not interchangeable. Where $\tau_{\mathrm{corr}}$ is shortest
--- at perilune, where the degree is highest --- the vehicle spends the least
time, and the integrator's call rate is highest there as well, so both factors
move together. Reading the local rate at perilune against an arc-averaged call
rate overstates the cost substantially.

Evaluated on the archived reference arcs, Eq.~\eqref{eq:probe-overhead-arc}
gives \ph{pct} for a near-circular low orbit and \ph{pct} for the eccentric
geometries of Table~\ref{tab:populations}, the eccentric cases being
\emph{cheaper} because their long apolune stretches need almost no probing.
Those estimates hold $r_{\mathrm{RHS}}$ fixed at its arc average and therefore
still err high; the measured values are in
Table~\ref{tab:probe-cost}. The probe remains the dominant term in $B_+$
--- against \ph{pct} for the pilot arc and \ph{pct} for switching --- which is
why it is reported as a curve rather than absorbed.

Two consequences follow, and neither is a footnote. The overhead is reported as
a \emph{measured curve} rather than a number: Table~\ref{tab:probe-cost} gives
the effective direction accuracy $\kappa_{\mathrm{eff}}$ obtained by probing at
$n_{\mathrm{probe}}$ points, averaged over the interval, against the cost that
buys it. And a cheap endpoint of that curve is declared in advance:
\Clite\ raises the degree of the interval's \emph{first} right-hand-side call
--- a point the integrator evaluates anyway, so $2k/N$ genuinely applies ---
and accepts that the direction is known only at the interval's start, for
$\tau_{\mathrm{corr}}$ of its length. \wpref{WP5, WP16}

\paragraph{Switching is not the expensive part.} A decision boundary makes the
right-hand side discontinuous, and one might expect that to dominate. It does
not. What the previous campaign measured is the aggregate: a radially binned
schedule cost \ph{pct} extra right-hand-side calls over seven days. Converting
that to a per-switch figure requires counting the bin crossings, which we do
from the archived degree histories rather than estimating; on \ph{n} crossings
it is \ph{value} extra calls per switch. At
$\Delta t_{\mathrm{dec}}=\ph{value}$~s there are \ph{n} boundaries per arc, so
the switching term is \ph{pct} --- roughly an order of magnitude below the
probe.

\begin{table}[!htbp]
  \centering\small
  \caption{The cost--accuracy curve of the online probe, measured rather than
  counted. $\kappa_{\mathrm{eff}}$ is the direction accuracy averaged over the
  decision interval when the interval is probed at $n_{\mathrm{probe}}$ points;
  the cost column is the share of $B_{\mathrm{tot}}$ the probe consumes, which
  Eq.~\eqref{eq:matching} subtracts from the controller's gravity budget. The
  first row is \Clite, whose single probe is co-located with a right-hand-side
  call and therefore costs an increment rather than a synthesis.
  \ph{fill from WP5, WP16}}
  \label{tab:probe-cost}
  \begin{tabular}{@{}lccccc@{}}
    \toprule
    Variant & $n_{\mathrm{probe}}$ & $\delta$ (idx) & $\Delta_{\mathrm{span}}$ (deg)
      & $\kappa_{\mathrm{eff}}$ & cost, share of $B_{\mathrm{tot}}$ \\
    \midrule
    \Clite & 1 (co-located) & \ph{n} & \ph{n} & \ph{value} & \ph{pct} \\
    \Cplan & \ph{n} & \ph{n} & \ph{n} & \ph{value} & \ph{pct} \\
    \Cplan & \ph{n} & \ph{n} & \ph{n} & \ph{value} & \ph{pct} \\
    \bottomrule
  \end{tabular}
\end{table}

\begin{figure}[!tbp]
  \centering
  \ph{figures/fig\_oa\_probe.pdf}
  \caption{Probe direction accuracy $\kappa=\cos\angle(\hat{\mathbf v},
  \Delta\mathbf a)$ against probe depth $k$, truncation degree $N$ and
  altitude. The dashed line is the level below which the coupling term cannot
  be carried (Section~\ref{sec:controller-limits}).}
  \label{fig:probe}
\end{figure}

\ph{one paragraph} Median $\kappa$ is \ph{value} at $k=1$ and \ph{value} at
$k=3$; it \ph{degrades/holds} with altitude and \ph{degrades/holds} with $N$.
The refresh interval required to keep $\kappa$ above \ph{value} is
\ph{value}~s, consistent with Eq.~\eqref{eq:decorrelation}.

\paragraph{What $\kappa$ does and does not measure.} It is tempting to read
$\kappa$ as the fraction of the coupling term the probe retains. It is not,
and the difference matters enough to state. The term is
$2\,\mathbf u_i(N)^{\top}\mathbf c_i$ with
$\mathbf u_i=\Phi\mathbf B\,\Delta\mathbf a\,\Delta t_i$, so writing
$\mathbf z=(\Phi\mathbf B)^{\top}\mathbf c_i$ the retained fraction is
$\langle\hat{\mathbf v},\mathbf z\rangle/
 \langle\Delta\mathbf a,\mathbf z\rangle$: it depends on the angles the true
and estimated defects make with $\mathbf z$, not on the angle between them,
because $\Phi\mathbf B$ does not preserve angles. Where
$\Delta\mathbf a\perp\mathbf z$ the true contribution vanishes and any probe
error produces a spurious one.

$\kappa=1$ still forces the retained fraction to one, and a low $\kappa$ still
destroys it, so $\kappa$ is a leading indicator rather than the quantity
itself --- and it has the practical advantage of being measurable without a
schedule and without propagating. We therefore report both: $\kappa$ as the
leading indicator, and the retained fraction, computed with $\mathbf c_i$ from
the \Asens\ schedule, alongside it. \emph{Both are diagnostics.} Neither
carries a stopping rule; the deployability decision is the capture fraction of
Section~\ref{sec:campaign}, which measures the thing itself rather than a
proxy for it. \wpref{WP5}

\subsection{Planning: forward--backward allocation}
\label{sec:controller-ifbda}

The coupling term $\mathbf c_i$ of Eq.~\eqref{eq:cross-term} depends on the
degrees chosen at \emph{every other} epoch. A single backward sweep of an
adjoint cannot supply it, because there is no schedule yet to sweep against.
The controller therefore iterates.

\begin{algorithm}[!htbp]
\caption{Forward--backward degree allocation on a pilot arc}
\label{alg:ifbda}
\begin{algorithmic}[1]
\Require pilot arc; low-degree $\Phi$; tail model; budget $B$
\State $N^{(0)} \gets$ separable sensitivity-weighted schedule (\Asens\ form)
\For{$j = 1 \dots J$}
  \State \textbf{forward:} accumulate $\hat{\mathbf u}_i$ from
    $\hat{\mathbf v}(t_i,N^{(j-1)}_i)$; form prefix sums
  \State \textbf{backward:} form the suffix blocks $\mathbf A_i$ and the cross
    terms $\mathbf c_i$
  \State \textbf{update:} $N^{(j)}$ by Eq.~\eqref{eq:coordinate-update};
    bisect $\lambda$ to the budget
\EndFor
\State \Return cancellation target $\mathbf c$, diagonal kernel $\mathbf K_i$,
  multiplier $\lambda$
\end{algorithmic}
\end{algorithm}

$J$ is a free parameter and is swept in Section~\ref{sec:ablations};
\ph{value} iterations suffice on \ph{n} of \ph{n} panel orbits.

\subsection{Probing forward, not backward}
\label{sec:controller-causality}

A decision interval is $\Delta t_{\mathrm{dec}}\approx\ph{value}$~s long and the
omitted direction decorrelates in $\Delta t_{\mathrm{probe}}\approx\ph{value}$~s.
A rule that accumulated probes taken \emph{inside} an interval and applied them
at the \emph{next} boundary would therefore select a degree from a direction
that has already decorrelated --- exactly the failure the resolution-scale
argument of Section~\ref{sec:controller-direction} predicts. The controller
does the opposite: at each boundary it probes \emph{ahead}, over the interval it
is about to commit to.

That is legitimate, and the reason is that two different coherence questions are
involved and they have different answers.

\begin{table}[!htbp]
  \centering\small
  \caption{Two coherence scales. The long-horizon failure is what forces the
  plan/probe split of Section~\ref{sec:controller-split}; the short-horizon
  success is what makes forward probing admissible.}
  \label{tab:coherence}
  \begin{tabular}{@{}p{0.24\textwidth}p{0.30\textwidth}p{0.36\textwidth}@{}}
    \toprule
    Horizon & State prediction error & Against $\pi r/N$ \\
    \midrule
    One decision interval,
    $\Delta t_{\mathrm{dec}}$ & the vehicle's own short-arc prediction,
      \ph{value}~m & far below; forward probing valid \\
    \addlinespace
    Whole arc, pilot against
    flown & \ph{value}~km after $T_{\mathrm{coh}}$ & far above; a pre-computed
      $\hat{\mathbf v}$ is invalid \\
    \bottomrule
  \end{tabular}
\end{table}

\paragraph{What predicts the probe locations.} At $t_q$ the degree for the
upcoming interval has not been chosen, so the integrator has not yet advanced
through it, and the probe points have to come from somewhere else. Two facts
make this a small requirement rather than a circularity.

First, the probe locations barely depend on the degree at all. Two candidate
degrees differ in acceleration by the difference of their truncation defects,
so over one interval they separate by at most
$\tfrac12\lVert\Delta\mathbf a\rVert\Delta t_{\mathrm{dec}}^2$, which for
$\lVert\Delta\mathbf a\rVert\sim\ph{value}$ and
$\Delta t_{\mathrm{dec}}=\ph{value}$~s is \ph{value}~mm. Against $\pi r/N$ that
is \ph{value} orders of magnitude below the scale on which the direction
changes. The predictor therefore does not need to know $N_q$, which is what
breaks the circle.

Second, what the predictor must achieve is a position accurate to well inside
$\pi r/N$ --- kilometres --- not an accurate trajectory. The controller uses a
two-body propagation of the current state over the interval, which requires no
field evaluation and is analytic. Its error is dominated by the perturbing
acceleration it omits, $\tfrac12\lVert\mathbf a_{\mathrm{pert}}\rVert
\Delta t_{\mathrm{dec}}^2\approx\ph{value}$~m, again far inside the tolerance.
Its cost is \ph{value} and enters $B_+$ with the rest.

\code{abl-predictor} replaces it with a low-degree micro-propagation at
$N_{\mathrm{plan}}$ and reports the change in $\kappa$ and in the final error,
so the choice is measured rather than assumed \wpref{WP5}.

The controller therefore evaluates the shared band stack at
$n_{\mathrm{probe}}=\lceil\Delta t_{\mathrm{dec}}/\Delta t_{\mathrm{probe}}\rceil$
points spanning the upcoming interval, along that predicted path, and sums
their transported contributions to obtain each candidate's
$\sum_{i\in I_q}\mathbf u_i(N)$. Nothing from before the boundary is reused for
the direction.

\begin{algorithm}[!htbp]
\caption{Online decision at boundary $t_q$}
\label{alg:online}
\begin{algorithmic}[1]
\Require plan $(\mathbf c,\mathbf K_i,\lambda,N_{\mathrm{plan}})$; window half-width
  $\delta$ (in candidate-grid indices); depth $k$; spent work $W$
\State predict the state at $n_{\mathrm{probe}}$ points spanning
  $[t_q,t_{q+1})$ by two-body propagation of the current state
  \Comment{no field evaluation; degree-independent}
\State evaluate bands $N_{\mathrm{plan}}(t_q)-\delta+1 \dots
  N_{\mathrm{plan}}(t_q)+\delta+k$ at those points \Comment{one shared stack}
\For{$N \in \mathcal{N}_q$}
  \State $\hat{\mathbf v}(t_i,N) \gets -\gamma(N,r_i)\sum_{n=N+1}^{N+k}
    \mathbf a_n(t_i)$ \Comment{partial sums, no new evaluation}
\EndFor
\State $N_q \gets$ Eq.~\eqref{eq:coordinate-update} over $\mathcal{N}_q$ with
  $\mathbf c$ from the plan
\State $\lambda \gets \lambda + \mathrm{gain}\cdot(W - W_{\mathrm{sched}}(t_q))$
  \Comment{budget feedback, Section~\ref{sec:controller-feedback}}
\State hold $N_q$ over $[t_q,t_{q+1})$
\end{algorithmic}
\end{algorithm}

Holding the degree across the interval also gives the controller exactly the
switching freedom the benchmark of Section~\ref{sec:problem} has, so the two
differ in information and not in resolution.

\dnote{If the decision grid has to shrink toward
$\Delta t_{\mathrm{probe}}$ for the coupling term to be worth anything, say so:
that is a finding about the controller's required cadence, and it should be
reported with the integrator cost it provokes rather than hidden by choosing
$\Delta t_{\mathrm{dec}}$ after the fact. The sweep in
Section~\ref{sec:ablations} is what decides this.}

\subsection{Why the plan and the probe are split}
\label{sec:controller-split}

Algorithm~\ref{alg:ifbda} runs on a pilot arc, which is not the arc that will
be flown. Whether that matters depends on which quantity is being carried
across, and the two behave differently.

The local omitted vector $\hat{\mathbf v}$ is field texture and decorrelates
over $\pi r/N$. A pilot arc's position error exceeds that scale after
$T_{\mathrm{coh}}=\ph{value}$, measured in Section~\ref{sec:controller-limits},
which is \ph{fraction} of the arc --- the long-horizon row of
Table~\ref{tab:coherence}. A schedule that fixed the probe at plan time would
therefore be using a direction that is no longer the right one, and the
controller does not do that: $\hat{\mathbf v}$ is evaluated online, at the
vehicle's own state, one interval at a time
(Section~\ref{sec:controller-causality}).

The cancellation target $\mathbf c$ is an accumulated quantity. The previous
paper measured the defect to be coherent rather than incoherent --- growth
exponents of \ph{range} against the $t^{3/2}$ of a zero-mean impulse null ---
so the accumulated displacement is dominated by a sustained trend rather than
by local texture, and it is stable between the pilot and the reference arc to
\ph{value} \wpref{WP6}. The same is true of the sensitivity weight and of the
multiplier.

The controller is accordingly split:

\begin{center}
\begin{tabular}{@{}ll@{}}
  \toprule
  Computed offline, on the pilot arc & $\mathbf c$, $\mathbf K_i$, $\lambda_0$,
    $N_{\mathrm{plan}}$, \emph{indexed by orbital phase} \\
  Computed online, ahead of each boundary & $\hat{\mathbf v}$ by
    Eq.~\eqref{eq:probe}, hence $N_q$ within $\pm\delta$ of the plan \\
  \bottomrule
\end{tabular}
\end{center}

and neither half reads the reference field or the reference arc.

\paragraph{The plan is indexed by phase, not by time.} This is forced by the
same coherence argument, applied to the smooth quantities rather than to the
texture. A loose $N=40$ pilot arc drifts along track by of order \ph{value}~km
over seven days, which at low-lunar-orbit speed is a \emph{timing} offset of
\ph{value}~s. A perilune passage lasts a few minutes. Since $\mathbf c$ and
$\mathbf K_i$ are dominated by exactly those passages, a plan indexed by
absolute time would
apply the perilune cancellation target before or after the perilune it was
computed for --- an error of order the thing being corrected.

Indexing by (revolution number, phase within revolution) removes the
sensitivity: the vehicle knows its own phase, and a \ph{value}~s offset is
\ph{value} of a revolution. It is not a return to a single-valued function of
radius, because the revolution index carries the remaining horizon; that is
what distinguishes the plan from the radial rules of
Table~\ref{tab:ladder}. \wpref{WP6, WP9}

\subsection{Budget feedback}
\label{sec:controller-feedback}

A per-call budget does not survive contact with an adaptive integrator: the
solver takes short steps where a concentrated schedule runs its highest
degrees, so nominal equality becomes a realized overspend --- a median
\ph{pct} in the previous campaign, and a factor of \ph{value} at half the
budget. \Cfb\ carries the spent work as a state, accumulating $N^2$ over the
right-hand-side calls actually made and adjusting $\lambda$ so that the
declared budget is met end to end rather than per call. \ph{one paragraph on
the update law and its gain.} The realized-to-nominal ratio it achieves is
\ph{value} against \ph{value} for the open-loop schedule.

\subsection{Where the controller is not the benchmark}
\label{sec:controller-limits}

Four separations are measured rather than argued, and each is an ablation in
Section~\ref{sec:ablations}: the tail amplitude in place of the true defect
magnitude; the probe direction in place of the true direction; the pilot arc in
place of the reference arc; and the low-degree $\Phi$ in place of the exact
one. \ph{one paragraph summarizing which dominates.}

\dnote{This is where a low $\kappa$ would first show. It is a warning and not a
stopping rule --- $\kappa$ is not the retained fraction, so gating on it would
risk an early false negative --- but if the capture fraction then confirms it,
the coupling term cannot be carried, \Cplan\ collapses onto rung 3, and the
paper's constructive claim reduces to a sensitivity-weighted controller ---
which is still new relative to the previous work, but is a different and
smaller claim. Draft that version of this section before the number arrives.}

\section{Propagated campaign}
\label{sec:campaign}

Sections~\ref{sec:benchmark} and~\ref{sec:controller} are first-order
statements about a fixed reference trajectory. This section propagates.

\subsection{Screening}

Each augmented variational solve costs about one propagation, and a propagated
comparison costs several, so candidates were screened before they were flown.
The screening instrument is the forced variational solution, which in the
previous campaign reproduced the sign of the propagated ratio on \ph{n} of
\ph{n} resolved comparisons of a \ph{n}-orbit panel. All candidates share one
reference trajectory and one gradient inside a single augmented integration and
differ only through their $\Delta\mathbf a$, so a panel of \ph{n} orbits is
solved once per orbit with every candidate as a separate forcing channel rather
than once per candidate. The selection rule, the panel and the number of
candidates allowed through --- two --- were fixed before the screen ran.
\ph{list the candidates and the outcome.} \wpref{WP13}

\paragraph{One caution about the screen, stated because it is load-bearing.}
The previous campaign's calibration of this instrument was established on
policies that were \emph{not} constructed by minimizing $J$. \Asign\ is, and so
are the controllers distilled from it. Scoring a candidate with the same
functional it was optimized against ranks candidates fairly but does not test
whether $J$ predicts $E$ for that class of schedule --- it assumes it. That
assumption is what the sign check below tests, which is why a failure there is
treated as a statement about the instrument rather than as noise, and why the
screening escape of Section~\ref{sec:setup-prereg} bounds how long it may be
diagnosed.

\subsection{What is propagated}

Five policies are propagated at each campaign point, and two of them are there
for reasons worth stating. \Asign\ is propagated because the first-order
benchmark has to be tested as a trajectory statement, not assumed to be one,
and because it is the denominator of the capture fraction below. \Rint\ is
propagated because it is the strongest deployable schedule available before
this paper, so beating \Fop\ while losing to \Rint\ would not be a result.

\begin{center}\small
\begin{tabular}{@{}ll@{}}
  \toprule
  Propagated & \Cplan, \Cfb\ (or the two screened candidates), \Asign, \Fop, \Rint \\
  Not propagated & \Fenv, constructed as a lower envelope over the constant
    family sweep \\
  \bottomrule
\end{tabular}
\end{center}

Each of the four calibrated policies --- the two screened controllers,
\Asign\ and \Rint\ --- is brought onto the anchor $B_{\mathrm{tot}}$ by the
iteration of Section~\ref{sec:setup-budgets}, which is why the campaign
propagates more arcs than it compares: \ph{value} propagations per policy per
orbit, \ph{n} in total. \Fop\ is the anchor and is not calibrated.
Table~\ref{tab:beta1} reports the achieved match alongside each verdict, and no
comparison is read where the match falls outside the declared band.

Arcs from the previous campaign are reused only where their realized work can
be matched under the same convention; where it cannot, the arc is propagated
again. Every reused arc is flagged as such in the record, because the earlier
campaign matched a different budget definition and a silent reuse would charge
that difference to the method. \ph{n} arcs were reused and \ph{n} were
recomputed.

\subsection{Equal realized work at $\beta=1$}

\begin{table}[!htbp]
  \centering\small
  \caption{The controller against both constant-degree comparators at
  $\beta=1$, matched on realized total quadratic work including the
  controller's own overhead $B_+$. Counts are per-orbit comparisons clearing
  the summed reference-inclusive envelope; $\rho$ is the median of per-orbit
  ratios over all orbits of the design, not over the resolved subset. Every
  candidate sits on the anchor $B_{\mathrm{tot}}$ to within \ph{pct}; the
  achieved match is archived per orbit and its worst case quoted here.
  \ph{fill from WP14}}
  \label{tab:beta1}
  \begin{tabular}{@{}lcccccc@{}}
    \toprule
    & \multicolumn{3}{c}{Design A} & \multicolumn{3}{c}{Design B} \\
    \cmidrule(lr){2-4}\cmidrule(lr){5-7}
    Comparator & wins & losses & $\rho$ & wins & losses & $\rho$ \\
    \midrule
    \Fop\ (deployable)   & \ph{n} & \ph{n} & \ph{value} & \ph{n} & \ph{n} & \ph{value} \\
    \Rint\ (deployable)  & \ph{n} & \ph{n} & \ph{value} & \ph{n} & \ph{n} & \ph{value} \\
    \Fenv\ (post hoc)    & \ph{n} & \ph{n} & \ph{value} & \ph{n} & \ph{n} & \ph{value} \\
    \bottomrule
  \end{tabular}
\end{table}

\ph{two paragraphs} The verdict against \Fop\ and \Rint\ is the claim; the row
against \Fenv\ is reported because a reader is entitled to know how much of the
constant family's potential the deployable comparator leaves unused, and it is
not a policy anyone could select in advance.

\begin{figure}[!tbp]
  \centering
  \ph{figures/fig\_oa\_beta1\_per\_orbit.pdf}
  \caption{Orbit-by-orbit ratio at $\beta=1$, sorted within each design.
  Colour records each orbit's own resolution verdict, so grey bars are
  comparisons the experiment cannot decide rather than ties.}
  \label{fig:beta1}
\end{figure}

\subsection{Capture fraction}

The controller is measured against the benchmark it approximates:

\begin{equation}
  f \;=\; \frac{E_{\Fop}-E_{\Cplan}}{E_{\Fop}-E_{\Asign}} ,
  \qquad\text{defined only where}\quad
  E_{\Asign} < E_{\Fop}
  \;\text{ and }\;
  M_{\mathrm{res}}(\Fop,\Asign) > 1 .
  \label{eq:capture}
\end{equation}

Three details are fixed in advance. The benchmark \Asign\ is not deployable and
carries no overhead, so it receives the full $B_{\mathrm{tot}}$ while \Cplan\ receives $B_{\mathrm{tot}}-B_+$; that asymmetry is deliberate, and it is why
$f$ measures the price of deployability rather than only the quality of the
approximation. Both terms are propagated errors $E$, never predictions, so $f$
does not exist before this section. And $f>1$ is possible without being an
error: \Asign\ is a first-order optimum on the reference arc and can be beaten
on the arc it is actually flown along, which is what
Section~\ref{sec:campaign-fixedpoint} tests separately.

The validity condition is not cosmetic. Where the benchmark does not beat the constant
degree the denominator is negative and $f$ has no interpretation; where the two
are separated by less than the numerical envelope it is a ratio of two
unresolved quantities and diverges. Orbits failing either test are reported as
\emph{not applicable} and counted, never coerced to zero or to a large number.
\ph{n} of \ph{n} orbits are eligible. \ph{one paragraph with the median $f$ by
design and budget.}

\begin{figure}[!tbp]
  \centering
  \ph{figures/fig\_oa\_capture.pdf}
  \caption{Capture fraction across the budget grid, both designs.}
  \label{fig:capture}
\end{figure}

\subsection{Across the budget grid}

\ph{one paragraph plus table} $\beta\in\{0.50,0.75,1.00,1.25,1.50\}$. The
compute-starved end is where every deployable alternative lost in the previous
campaign, so it is reported first and separately. \ph{state the verdict at
$\beta=0.50$ explicitly, whichever way it goes.}

\dnote{H4 and H5 are decided here and are separate hypotheses on purpose: H5
($\beta=0.5$) can be rejected without touching H4. Report a rejection as a
rejection.}

\subsection{Across populations}

\ph{one paragraph plus table} Designs A, B, C; the five geometry strata; the
wide-elliptic population, which is the regime where the previous paper's radial
endpoint \emph{won}, so a trajectory-aware controller that fails there is
regime-bound; and the low-perilune subset, where the linearization behind both
the benchmark and the planner is weakest and only signs are read.

\subsection{Cost ladder}

\begin{table}[!htbp]
  \centering\small
  \caption{What ``equal budget'' equals, down the four costs of
  Section~\ref{sec:setup-budgets}, as controller over comparator. Only the
  primary row is matched by construction. \ph{fill from WP16}}
  \label{tab:cost-ladder}
  \begin{tabular}{@{}lccl@{}}
    \toprule
    Cost & A & B & \\
    \midrule
    Realized total quadratic work $B_2$ & \ph{value} & \ph{value} & matched by construction \\
    \quad of which controller overhead $B_+$ & \ph{pct} & \ph{pct} & pilot arc, probe, planning \\
    Nominal per-call work $B_1$ & \ph{value} & \ph{value} & reported for comparability only \\
    Right-hand-side calls made & \ph{value} & \ph{value} & \\
    End-to-end wall time & \multicolumn{2}{c}{\ph{value}} & on \ph{n} timed orbits \\
    Measured kernel time $B_3$ & \multicolumn{2}{c}{\ph{value}} & idle machine, three repeats \\
    \bottomrule
  \end{tabular}
\end{table}

\ph{one paragraph} The overhead row is the one to read first: a controller that
wins on error while spending its margin on its own planning has not won.

\subsection{Does the first-order benchmark survive propagation?}
\label{sec:campaign-fixedpoint}

\Asign\ optimizes a linearization about a fixed reference arc. When its
schedule is flown, the trajectory changes, $\Phi$ changes, and the local defect
changes with them, so the schedule that was optimal for the reference arc need
not be optimal for its own. The fixed-point test propagates the schedule,
re-tabulates along the resulting arc, re-solves, and propagates again.
\ph{one paragraph: how far the schedule moves, how far the error moves, and
after how many iterations.} \wpref{WP17}

\dnote{This is the test that decides whether \Asign\ is a
\emph{trajectory-level} benchmark or a \emph{reference-arc} one. If the gain
does not survive, say so in these words and rename the quantity throughout;
that outcome is a strong negative result about signed allocation and is worth
stating cleanly rather than burying.}

\subsection{Controls}

Each was drawn under a rule fixed before the runs it selects, and each targets
a way the result could be an artefact rather than a measurement.

\begin{table}[!htbp]
  \centering\small
  \caption{Registered controls. \ph{fill}}
  \label{tab:controls}
  \begin{tabular}{@{}p{0.24\textwidth}p{0.30\textwidth}p{0.38\textwidth}@{}}
    \toprule
    Control & What it tests & Outcome \\
    \midrule
    Reference degree $300\to600$ & whether large ratios are
      reference-adjacency artefacts rather than measurements & \ph{outcome} \\
    Decision-grid parity & that \Asign\ and \Cplan\ switch on the same grid, so
      the benchmark's margin is information and not resolution & \ph{outcome} \\
    Degree-cap audit & orbits whose schedule reaches the reference degree,
      where its force model \emph{is} the reference model & \ph{outcome} \\
    Phase-shifted initial states & whether the signed allocation is fitted to
      one arc's cancellation pattern & \ph{outcome} \\
    Time-indexed against altitude-binned schedule & whether the gain belongs to
      the schedule's form & \ph{outcome} \\
    Decision-grid coarsening & whether it belongs to the decision cadence &
      \ph{outcome} \\
    Information-source audit & that no \Cplan\ run read the reference field or
      the reference arc & \ph{outcome} \\
    \bottomrule
  \end{tabular}
\end{table}

\section{Ablations: which ingredient earns the gain}
\label{sec:ablations}

Each ablation removes one ingredient and leaves the rest intact, and all are
scored by the forced variational solution rather than by propagation.

Two things about the panel are declared rather than left to be inferred. The
fourteen ablations are run on a \ph{n}-orbit perilune-spread panel, smaller
than the \ph{n}-orbit panel the candidates themselves were screened on, because
each variational solve costs about one propagation and fourteen removals across
the full panel would cost more than the propagated campaign it is meant to
inform. And all fourteen share one augmented integration per orbit, differing
only in their forcing channel, so the panel is the same for every row of
Table~\ref{tab:ablations} and the rows are comparable with each other. They are
\emph{not} directly comparable with the ladder of Table~\ref{tab:ladder}, which
stands on the full design. \wpref{WP9, WP13}

\begin{table}[!htbp]
  \centering\small
  \caption{Ablations at $\beta=1$. ``Retained'' is the fraction of the
  controller's advantage over \Fop\ that survives the removal. \ph{fill}}
  \label{tab:ablations}
  \begin{tabular}{@{}llcc@{}}
    \toprule
    Code & What is removed & Retained & Note \\
    \midrule
    \code{abl-sign}  & the coupling term; \Cplan\ collapses to a
      sensitivity-weighted rule & \ph{pct} & answers RQ2 on the deployable side \\
    \code{abl-phi} & $\Phi$, leaving only a $(T-t)$ horizon weight
      (\Ctgo) & \ph{pct} & \\
    \code{abl-probe} & the probe; direction replaced by a fitted
      $(r,\phi,\lambda,N)$ surrogate & \ph{pct} & tests
      Section~\ref{sec:controller-direction} directly \\
    \code{abl-null}  & the probe direction randomized at matched magnitude &
      \ph{pct} & negative control; must destroy the gain \\
    \code{abl-kaula} & the measured spectrum $P_n$ in $\gamma$ replaced by a
      fitted power law & \ph{pct} & prices the choice of
      Eq.~\eqref{eq:gamma} over a Kaula extrapolation \\
    \code{abl-pilot} & pilot arc replaced by the reference arc & \ph{pct} &
      prices the foreknowledge \\
    \code{abl-stm}   & low-degree $\Phi$ replaced by a Keplerian one &
      \ph{pct} & is a closed form enough \\
    \code{abl-J}     & $J=1$ instead of $J=\ph{value}$ & \ph{pct} & how many
      passes are needed \\
    \code{abl-k}     & probe depth $k=1$ instead of $k=\ph{value}$ &
      \ph{pct} & \\
    \code{abl-window} & candidate window $\delta$ narrowed and widened &
      \ph{pct} & prices the plan/probe split; wider costs $B_+$ \\
    \code{abl-grid}  & decision grid $\Delta t_{\mathrm{dec}}$ swept over
      \ph{value}~s & \ph{pct} & the cadence the coupling needs, against the
      integrator work it provokes \\
    \code{abl-lookahead} & forward probe replaced by a probe at the boundary
      only & \ph{pct} & tests Section~\ref{sec:controller-causality} directly \\
    \code{abl-predictor} & two-body probe-point predictor replaced by a
      low-degree micro-propagation & \ph{pct} & prices the predictor choice \\
    \code{abl-timeindex} & phase-indexed plan replaced by a time-indexed one &
      \ph{pct} & prices the pilot arc's timing offset,
      Section~\ref{sec:controller-split} \\
    \bottomrule
  \end{tabular}
\end{table}

\begin{figure}[!tbp]
  \centering
  \ph{figures/fig\_oa\_ablation\_waterfall.pdf}
  \caption{Ablation waterfall from the benchmark down to the constant degree.}
  \label{fig:waterfall}
\end{figure}

\paragraph{The negative control is the one to read first.} \code{abl-null}
keeps every magnitude and randomizes only the direction. If the controller
retains its advantage under that ablation, the advantage is not coming from the
signed information the paper claims it comes from, and no other row means what
it appears to mean. \ph{state the outcome.}

\paragraph{The surrogate ablation tests an argument, not a variant.}
Section~\ref{sec:controller-direction} argues that tabulating the omitted
direction is not a cheap alternative, because it varies on the field's own
resolution scale.
\code{abl-probe} builds the surrogate anyway --- \ph{describe: fitted on the
reference field over the sampled altitude range, with \ph{n} degrees of
freedom} --- and measures what it retains. \ph{one paragraph.} What is
reported is the pair: what the surrogate retains, and what it costs in degrees
of freedom and evaluation time to retain it. A surrogate that resolves
$\pi r/N$ carries a number of parameters comparable to the omitted
coefficients, \ph{value} against \ph{value}, and that comparison is the
measured form of the argument in Section~\ref{sec:controller-direction} --- not
a claim that no such surrogate can be built.

\subsection{Four controls on the design of the controller itself}
\label{sec:ablations-design}

The ablations above price the ingredients of \Cplan. They cannot say whether
\Cplan\ is the right \emph{shape}, because every row holds the architecture
fixed. Four further controls test the architecture. None of them requires a
propagation: all four run on the tabulated $\mathbf u_i(N)$, the suffix
sequence $\mathbf A_i$ and the Frank--Wolfe iterate that
Sections~\ref{sec:benchmark} and~\ref{sec:algorithm} produce anyway.
\wpref{WP21}

\paragraph{Is the descent the weak link, or is the bound?} The certified gap
$g_E$ conflates two quantities --- how far the descent sits above the integer
optimum, and how far the relaxation sits below it --- and the campaign's
response to a wide gap differs completely between them. On short instances
($K_{\mathrm{dec}}=\ph{n}$, $\lvert\mathcal N\rvert=\ph{n}$) the integer
optimum $J^{*}$ is obtained exactly by enumeration, which splits the gap into
$J_{\mathrm{desc}}-J^{*}$ and $J^{*}-L_{\mathrm{FW}}$. \ph{state which
dominates.} We also report a second solver on the full problem:
\code{S-round} rounds the Frank--Wolfe iterate --- per-interval
$\arg\max_N\theta_{qN}$,
then sum-up rounding, then one polishing sweep --- and if it beats the
multi-start descent it becomes the reported schedule and the descent becomes
the control.

\paragraph{How many directions does the coupling actually use?} Each
$\mathbf M_j$ has rank at most three, but $\mathbf A_i$ is a \emph{sum} of
them, and because the summands' ranges rotate along the arc the sum can have
rank up to six. Low rank is therefore something to measure, not something the
construction supplies.

\emph{Measured on a nondimensionalized form.} The state is
$\mathbf x=[\mathbf r^{\top}\ \mathbf v^{\top}]^{\top}$ with metres in the
first three components and metres per second in the last three, so the blocks
of $\mathbf A_i$ carry different physical dimensions and its raw eigenvalues,
trace and eigenvectors are not invariant under a change of units: reporting
velocity in km/s instead of m/s can change the effective rank. With
$\mathbf D=\operatorname{diag}(L,L,L,L/T_s,L/T_s,L/T_s)$ and
$\mathbf x=\mathbf D\tilde{\mathbf x}$, the quadratic form transforms
congruently,

\begin{equation}
  \tilde{\mathbf A}_i \;=\; \mathbf D^{\top}\mathbf A_i\,\mathbf D ,
  \label{eq:nondim-form}
\end{equation}

and every quantity below is computed on $\tilde{\mathbf A}_i$. This is a
congruence and not a similarity transform: it is the form that is being
rescaled, not an operator, and using $\mathbf D^{-1}\mathbf A_i\mathbf D$
would change the spectrum for no physical reason. The scales $L$ and $T_s$ are
registered in advance --- the lunar reference radius and the characteristic
orbital time of Table~\ref{tab:populations} --- so that the verdict cannot be
moved by choosing them afterwards.

\emph{T3a, concentration.} Table \ph{tab} gives $p$, the number of
eigendirections of $\tilde{\mathbf A}_i$ carrying $95\%$ of its trace, and the
objective error and schedule change when $\tilde{\mathbf A}_i$ is truncated to
them. \ph{state the number.}

\emph{T3b, coherence.} Concentration alone does not license a reduced
controller. If the dominant subspace $\mathcal U_p(t)$ turns quickly, a
$p$-scalar state is the wrong $p$ scalars one interval later. We therefore
report the principal angles between $\mathcal U_p(t_i)$ and
$\mathcal U_p(t_{i+1})$, and their accumulation over the coherence horizon of
Section~\ref{sec:controller-split}. \ph{state the median angle.}

The branch is taken on both together: $p\le2$ \emph{and} a median principal
angle below \ph{value} sends the architecture to a $p$-scalar state updated in
flight; either alone leaves the full six-vector target read from an offline
plan. The consequence is architectural rather than cosmetic, which is why the
rule is fixed here rather than after the figure exists.

\paragraph{How far ahead does the coupling reach?}
Table~\ref{tab:horizon} replaces $\mathbf A_i$ by its truncation to a horizon
$H$ and re-solves, which is \Asign\ run receding-horizon with perfect
information. The curve says what a replanning controller would need to look at,
and it is also the direct physical statement --- absent from
Section~\ref{sec:benchmark}, where it belongs --- of the range over which
cancellation operates.

\begin{table}[!htbp]
  \centering\small
  \caption{Horizon sufficiency. $\hat\rho(H)$ is the benchmark ratio when the
  allocation is solved against only the next $H$ of trajectory. $H^{*}$ is the
  smallest $H$ reaching $0.90$ of the full-horizon ratio. \ph{fill from WP21}}
  \label{tab:horizon}
  \begin{tabular}{@{}lcccccc@{}}
    \toprule
    $H$ (revolutions) & $1/8$ & $1/4$ & $1/2$ & $1$ & $2$ & full arc \\
    \midrule
    $\hat\rho(H)/\hat\rho(T)$ & \ph{v} & \ph{v} & \ph{v} & \ph{v} & \ph{v}
      & $1$ \\
    \bottomrule
  \end{tabular}
\end{table}

\paragraph{Could a state-feedback rule have found this without probing?} This
is the control that can falsify the paper's information claim, so it is stated
with its failure condition first. A policy is fitted to the \Asign\ schedules
on design A and evaluated on design B, and what is reported is the fraction of
\Asign's advantage over \Fop\ that it retains out of sample. In the smooth-state
variant the features are altitude, radius, speed, orbital phase, time to go and
budget spent --- position on the sphere is withheld. \textbf{If that variant
retains more than $0.60$ of the advantage, the gain is geometric rather than
textural, the probe of Section~\ref{sec:controller} is unnecessary, and this
paper's information statement is wrong.} \ph{state the outcome, and state it
as a rejection if it is one.} The second variant adds position at increasing
resolution and joins the surrogate curve above. The fitted policy is a control,
not a proposed method: the simplest expressive form is used and its complexity
is reported alongside what it buys.

\section{Discussion}
\label{sec:discussion}

\subsection{What a non-separable objective changes in practice}

\ph{one or two paragraphs} The separable reading of degree selection --- spend
where the local error is largest --- is not a poor approximation to the right
answer; it optimizes a different functional, and the previous paper measured
that the two can order two policies oppositely. What this paper adds is that
the right functional is still tractable, because the coupling has a prefix--
suffix structure, and that the part a deployable rule cannot easily obtain is
not the weight but the sign. The sharpest way to say it is that all three
criteria are readings of one trajectory kernel: the force-level rule ignores
it, the sensitivity-weighted repair retains only its local diagonal
$\mathbf K(t,t)$, and only the third retains the cross-epoch coupling that
encodes cancellation. The weight was never the missing ingredient, because the
weight \emph{is} that diagonal.

\subsection{The resolution scale as a design constraint}

\ph{one paragraph} One length scale, $\pi r/N$, turns out to fix five otherwise
unrelated design choices:

\begin{enumerate}
  \item the sampling requirement on the accumulation grid, so that the signed
    integral is not aliased (Eq.~\eqref{eq:decorrelation});
  \item the resolution a surrogate for the omitted direction would have to
    reach, and therefore its cost (Section~\ref{sec:controller-direction});
  \item the probe refresh interval;
  \item the probe's share of the budget, which is the arc integral of its
    reciprocal (Eq.~\eqref{eq:probe-overhead-arc}); and
  \item the coherence horizon of the pilot arc, which is what forces the plan
    to be phase-indexed and the probe to be online.
\end{enumerate}

The rule of thumb worth carrying away is
\emph{nothing that varies faster than the truncation's own resolution scale can
be planned in advance} --- and the corollary, which is the expensive half:
whatever must therefore be measured in flight costs a synthesis every time the
scale turns over.

\subsection{Relation to goal-oriented adaptivity}

\ph{one paragraph} \citep{becker2001optimal}. The structural difference is that
the functional here is quadratic in an integral of the residual, so the
adjoint weight is not a scalar field and the resulting allocation is coupled
rather than pointwise.

\subsection{When would a practitioner use this?}

\ph{one paragraph} The controller costs a pilot arc, an offline planning pass,
and an online probe of \ph{pct} of the gravity budget. Two of those three
amortize and one does not, which decides the answer. The pilot arc and the plan
are paid once per trajectory family and can be reused across a Monte Carlo
ensemble or a covariance study; the probe is paid on every arc, every time.
A single arc therefore pays the full setup for one use and is the worst case;
\ph{quantify the break-even in arcs.} If the campaign lands in the regime where
the probe alone exceeds the margin, the honest recommendation is \Clite\ or the
constant degree, and Section~\ref{sec:campaign} says which.

\subsection{Limits}

\begin{itemize}
  \item \textbf{Model-relative.} Everything is measured against an adopted
    reference expansion, not against the true lunar field.
  \item \textbf{Linearization.} Both the benchmark and the planner are
    first-order. Amplitudes are not interpreted below \ph{km} perilune, where
    the previous campaign found the first-order amplitude wrong by up to a
    factor of two while the sign survived.
  \item \textbf{Arc length.} Calibrated and evaluated on seven-day arcs.
    \ph{state the sixty-day outcome.} A schedule planned for one arc length
    does not transfer to another, because the perilune descends and the dwell
    shifts.
  \item \textbf{Geometry.} \ph{state which populations the claim holds on.}
  \item \textbf{Certificate scope.} The gap of
    Section~\ref{sec:algorithm-certificate} bounds the distance to the best
    schedule of the linearized reference-arc problem, not of the propagated
    one.
  \item \textbf{Overhead is not small.} The online probe consumes \ph{pct} of
    the budget, an order of magnitude more than the switching cost and enough
    to decide the comparison on its own. It is charged in full
    (Eq.~\eqref{eq:matching}), but a reader should treat any margin narrower
    than it as contingent on the measured overhead rather than on the
    allocation idea.
  \item \textbf{Screening is not independent of the objective.} Candidates were
    screened by the same functional the benchmark optimizes
    (Section~\ref{sec:campaign}). The propagated sign check is what tests that
    transfer, and it is load-bearing rather than confirmatory.
  \item \textbf{No located optimum.} No continuous optimum in $\beta$, in $k$,
    in $\delta$, in $n_s$ or in $J$ is located anywhere in this paper; all are
    sampled on declared grids.
  \item \textbf{The benchmark and the controller are not claimed with equal
    force.} \Asign\ is derived from the trajectory-error functional itself and
    is certified against it, so the attainability result stands on the
    derivation. \Cplan\ is not claimed to be the best deployable schedule ---
    it is the first interpretable, auditable, low-complexity one. Three of its
    design choices are measured rather than argued and each has a declared
    alternative: the coordinate descent against a global solve on short
    instances (Section~\ref{sec:ablations}), the frozen offline plan against
    the horizon over which the coupling actually reaches
    (Table~\ref{tab:horizon}), and the band probe against a state-feedback
    policy fitted to \Asign\ itself. Where a measurement favours the
    alternative, it is reported as favouring the alternative.
\end{itemize}

\subsection{Future work}

\ph{short paragraph} Each item below is named with the measurement in this
paper that decides whether it is worth doing, so the list is a set of
conditional branches rather than a wish list.

\begin{itemize}
  \item \emph{Receding-horizon replanning}, if the horizon curve of
    Table~\ref{tab:horizon} has its knee well inside the arc or the coherence
    horizon of Section~\ref{sec:controller-split} is short.
  \item \emph{A reduced cancellation state.} The coupling enters only through
    $\mathbf A_i$, whose rank is at most three by construction and whose
    eigenvalue spread is reported in Section~\ref{sec:ablations}. If the trace
    concentrates in one or two directions, the accumulated state the controller
    must carry is that many scalars rather than a six-vector, and it can be
    updated in flight instead of being planned offline --- a cheaper controller
    than the one flown here, and the same reduction that makes a trellis
    formulation tractable.
  \item \emph{A switching penalty}, which turns the coordinate descent into a
    trellis problem and is only worth the state if that reduction holds.
  \item \emph{A vector-tail surrogate} in place of the probe, along the
    cost--retention curve of Section~\ref{sec:ablations}, if the measured probe
    overhead stays at the level reported here.
  \item \emph{A stronger solver} --- relax-and-round on the Frank--Wolfe
    iterate, or a global solve --- if the gap decomposition attributes the
    certified gap to the descent rather than to the relaxation.
  \item Allocation against an ensemble of initial states rather than one;
    transfer to other fields \citep{konopliv2016mars}; and propagating
    coefficient uncertainty into the degree-variance spectrum the amplitude
    estimate rests on.
\end{itemize}

\section{Conclusion}
\label{sec:conclusion}

\dnote{Three versions of this section are drafted below, one per outcome band
of \code{../OUTCOMES.md}. Keep exactly one before submission and delete the
other two together with this note. Drafting them in advance is deliberate: the
point of the pre-registration is that the wording does not get chosen after the
numbers arrive.}

% ===========================================================================
% VERSION A -- constructive (green / amber-green band)
% ===========================================================================

Allocating a spherical-harmonic degree budget from the local force error
optimizes the wrong functional. Written in the form the dynamics impose, the
objective takes the norm after a signed integral against the state-transition
matrix, which couples every pair of epochs and makes the allocation
non-separable. The classical multiplier treatment and its sensitivity-weighted
repair both still \emph{work} --- they are separable and a single multiplier
decouples them --- but they optimize the wrong functional, and the discretized
form says exactly how: all three criteria are readings of one trajectory
kernel, which the force-level rule ignores, the sensitivity-weighted repair
reduces to its local diagonal, and only the third keeps whole. The weight was
never the missing ingredient, because the weight \emph{is} that diagonal.

The coupling nonetheless has a prefix--suffix structure that keeps a full sweep
over the decision intervals at $\mathcal{O}(M\lvert\mathcal{N}\rvert)$ and
admits a Frank--Wolfe bound on its convex-hull relaxation, so the problem is
solved and certified without propagating anything, at a certified error gap of
\ph{pct}. On \ph{n} orbits of \ph{n} independent designs the certified
allocation is a median \ph{value} better than the equal-budget constant degree
and \ph{value} better than the strongest previously available deployable
schedule, and \ph{pct} of that advantage lives off the diagonal.

Carrying it into a deployable rule turns on a quantity a tail criterion does
not provide: the \emph{direction} of the omitted acceleration. That direction
varies on the field's own resolution scale $\pi r/N$, which is \ph{value}~km at
degree 300, so a surrogate coarser than that scale cannot represent it and one
fine enough to do so carries a number of degrees of freedom comparable to the
omitted coefficients themselves. It can be probed, and the probe is the
controller's dominant cost rather than a rounding term: covering a decision
interval takes a measured \ph{pct} of the total budget, and recovers
the direction to $\kappa=\ph{value}$. The same length scale sets how far ahead
the probe may look and how long a plan made on a cheap pilot arc stays valid,
and those two answers differ: forward over one decision interval, yes; across a
seven-day arc, no. That is why the controller plans its cancellation target
offline and probes forward one interval at a time rather than uploading a
schedule.

Under a common realized-work ceiling, with its own overhead charged in full,
that controller has the smaller seven-day position error than the constant degree on
\ph{n} of \ph{n} resolved comparisons across \ph{n} populations, capturing a
median \ph{value} of the certified benchmark's advantage. \ph{state the
$\beta=0.5$ outcome explicitly.} \ph{state the wide-elliptic outcome
explicitly.}

Every conclusion is model-relative and bounded by the tested designs, arc
lengths and budgets. \ph{closing sentence naming the sharpest bound.}

% ===========================================================================
% VERSION B -- benchmark real, controller cannot capture it (amber band)
% ===========================================================================
%
% \ph{Draft:} the formulation and the certificate stand; the benchmark shows a
% median \ph{value} of attainable gain over the constant degree; no deployable
% controller captures a resolvable fraction of it. The decorrelation argument
% then carries the paper: what separates the benchmark from every deployable
% rule is not the weighting and not the algorithm but access to the omitted
% direction at the field's own resolution scale, and the probe measures exactly
% how much of that access is affordable. State the capture fraction, name which
% substitution destroys it (tail amplitude, pilot arc, low-degree $\Phi$, or
% probe depth), and present the result as an information-limit statement that
% strengthens rather than contradicts the previous paper.
%
% ===========================================================================
% VERSION C -- no attainable gain (orange band)
% ===========================================================================
%
% \ph{Draft:} even with the sensitivity and the sign, and with a certificate
% bounding how much better any schedule of this class could be, allocation does
% not systematically improve on a constant degree at equal realized work in
% this regime. That closes the question the previous paper left open rather
% than leaving it open, and it is worth stating without hedging. Distinguish
% the two ways it can happen --- no gain in the linearized problem at all, or a
% linearized gain that does not survive propagation
% (Section~\ref{sec:campaign-fixedpoint}) --- because they mean different
% things: the first says the objective has no exploitable structure at these
% budgets, the second says the structure is an artefact of optimizing along a
% fixed reference arc.


\section*{Acknowledgments}
\ph{to be written}

\section*{Data availability}
\ph{archive DOI, manifest digest chain, pre-registration hashes}

\bibliographystyle{plainnat}
\bibliography{references}

\end{document}
